\documentclass{article}
\usepackage[margin=0.8in]{geometry}
\usepackage{amsmath,amssymb,amsthm,mathtools}
\usepackage{mathrsfs,bm,bbm}
\usepackage{graphicx,booktabs,array,multirow,tabularx,longtable}
\usepackage{enumitem}
\usepackage{xcolor}
\usepackage{algorithm,algpseudocode}
\usepackage{subcaption,tikz}
\usepackage{siunitx,cancel,accents}
\usepackage{microtype}
\usepackage[numbers,sort&compress]{natbib}
\usepackage[colorlinks=true,linkcolor=blue,citecolor=blue,urlcolor=blue]{hyperref}
\usepackage{cleveref}
\setlength{\emergencystretch}{2em}
\newtheorem{assumption}{Assumption}
\newtheorem{definition}{Definition}
\newtheorem{theorem}{Theorem}
\newtheorem{lemma}{Lemma}
\newtheorem{proposition}{Proposition}
\newtheorem{openendedquestion}{Open-ended Question}
\newtheorem{conjecture}{Conjecture}
\newcommand{\coreref}[1]{\ref{#1}}

\begin{document}

\sloppy

\section*{stat\_\allowbreak{}unknowncutoff\_\allowbreak{}weaklate\_\allowbreak{}frontier}

\noindent\textit{Specialization:} v1\par

\paragraph{TL;DR.} The paper studies fixed-\(\alpha\) expected length of connected honest intervals for fuzzy-RD LATE when both the compliance jump and its unknown cutoff become statistically weak.  A parameter-exact outcome-invariant mixture gives the finite-sample lower bound with coefficient \(2(1-v)\) and a nonshrinking limit below the explicit witness threshold \(\lambda_w=315(M/32)^{1/s}/(256s)\).  This is a one-sided elbow, not the exact critical frontier.  At \(s=1\), a total four-fold rectangular scan--refine--invert rule attains the known-cutoff order when \(n\kappa_n^3\) dominates \(\log(e/\kappa_n)\).  For \(s>1\), the frozen affine moving-window profile is disproved as a uniform localizer, so the detectable-region bracket and the fixed-level critical-window law remain open.

\section{Project justification}

\paragraph{Gap.} Known-cutoff weak fuzzy-RD inference and fixed-jump unknown-cutoff localization do not characterize honest connected LATE-interval length when the first-stage denominator and its location weaken jointly.  Existing shrinking-jump location results do not yield the explicit causal witness elbow \(\lambda_w\), its parameter-exact \(2(1-v)\) connected-length coefficient, or an exact critical-window law.  In the higher-smoothness problem, general stochastic maximal inequalities are also insufficient because the frozen moving-window profile can be misranked already at the population level under separate one-sided smoothness.

\paragraph{Niche.} The joint experiment combines a nonconvex union over cutoff locations, a causal outcome--treatment ratio, separate one-sided observable smoothness that permits a derivative kink, and fixed-level connected-length loss.  The parameter-exact common-outcome construction makes the full likelihood depend only on \((X,D)\), so outcome-assisted localization cannot evade its explicit location entropy.  The higher-smoothness counterexample separately shows that a locally affine compliance trend can make a displaced profile split fit better than the true split.

\paragraph{Fill.} On one unchanged observable-law class, the project proves an explicit finite-sample location-mixture converse with coefficient \(2(1-v)\), evaluates its one-sided threshold \(\lambda_w=(sC_\psi C_h)^{-1}\), and gives a total four-fold interval attaining known-cutoff order on the Lipschitz face.  For \(s>1\), it proves an exact projection counterexample to the frozen profile and retains the proposed detectable-region bracket as an OEQ requiring a genuinely different estimator or profile, or a negative resolution.  The exact critical-window frontier is not claimed.

\section{Related work}

Noack and Rothe (2024) supply the known-cutoff bias-aware Anderson--Rubin template and separate smoothness restrictions on observable outcome and treatment regressions. Our binary-outcome, uniform-design, unknown-cutoff triangular class adapts that template but is not their published parameter space.  Porter and Yu (2015) obtain unknown-cutoff oracle substitution under a fixed discontinuity in a sharp-design treatment-effect problem.  Frick, Munk, and Sieling (2014) motivate multiscale scan architecture, while Chernozhukov, Chetverikov, and Kato (2014) provide general stochastic maximal-control tools; neither prevents the deterministic population-score misranking in the higher-smoothness counterexample.  M{\"u}ller and Song (1997) and Tuvaandorj (2020) concern shrinking-jump or nonparametric location experiments; neither converts translated-cutoff multiplicity into the proved \(2(1-v)\) fixed-level connected LATE-length bound below \(\lambda_w\), and no equation-level proposition from Tuvaandorj is used here.  Armstrong and Koles{\'a}r (2018) supply an ordered-modulus and test-inversion template for convex fixed-index classes.  Madrid Padilla, Yu, Wang, and Rinaldo (2022) study aligned multivariate change channels without the monotone fuzzy-RD causal restrictions or the common-outcome likelihood cancellation used here.  Kowalska, van de Wiel, and van der Pas (2025 version) use a Bayesian unknown-cutoff fuzzy-RD model, and Arai, Otsu, and Seo (2026) discuss fuzzy and unknown-discontinuity extensions separately.  The proved frontier advance is the explicit one-sided location-mixture converse together with Lipschitz oracle attainment; a replacement higher-smoothness procedure and the exact critical-window connected-length law remain open.

\section{Setup}

Target (estimand / structural parameter): \(\theta_P=y_P/d_P=E_P[Y(1)-Y(0)\mid R=C,X=c_P]\).
Primary functional / estimator / diagnostic: \(I_n^{\mathrm{SI}}=I_n^{\mathrm{SI},1}\) when \(s=1\); when \(1<s\le2\), it is the full-parameter fallback if \(\mathsf F_{\mathrm{prof}}=1\) or a displayed deterministic/denominator check fails, and otherwise it is \(\operatorname{hull}\{\theta\in\Theta:|\widehat y-\theta\widehat d|\le2r_{\mathrm{prof}}\}\) at the returned center \(\widehat c_{\mathrm{prof}}\)..
\begin{itemize}
  \item \texttt{n} --- \texttt{sample-\allowbreak{}size index}; \(\mathbb N\); \texttt{indexes the triangular experiment}
  \item \texttt{s} --- \texttt{smoothness exponent}; \([1,2]\); \texttt{governs the local regression class}
  \item \texttt{M} --- \texttt{smoothness radius}; \([2,\infty)\); \texttt{bounds observable one-\allowbreak{}sided reduced-\allowbreak{}form extensions}
  \item \texttt{v} --- \texttt{observable outcome-\allowbreak{}variance floor and correlation-\allowbreak{}separation constant}; \((0,1/64)\); \texttt{keeps transformed-\allowbreak{}outcome noise nondegenerate while permitting rare treatment}
  \item \texttt{\textbackslash{}\allowbreak{}alpha} --- \texttt{miscoverage level}; \((0,1/4)\); \texttt{fixes the honesty requirement}
  \item \texttt{\textbackslash{}\allowbreak{}lambda} --- \texttt{critical-\allowbreak{}window coordinate}; \((0,\infty)\); \texttt{indexes limits with signal information comparable to location entropy}
  \item \texttt{\textbackslash{}\allowbreak{}kappa\_\allowbreak{}n} --- \texttt{positive compliance-\allowbreak{}strength sequence}; \((0,1/8]\); \texttt{allows the fuzzy first-\allowbreak{}stage jump to vanish}
  \item \texttt{X} --- \texttt{running variable}; \([0,1]\); \texttt{indexes the unknown assignment threshold}
  \item \texttt{R} --- \texttt{monotone treatment-\allowbreak{}response type}; \(\{N,C,A\}\); \texttt{records never-\allowbreak{}taker,\allowbreak{} complier,\allowbreak{} or always-\allowbreak{}taker status}
  \item \texttt{D(z)} --- \texttt{binary potential treatment receipt}; \(\{0,1\}\); \texttt{encodes monotone compliance}
  \par\smallskip\noindent\emph{Definition.} \(D(z)=0\) for \(R=N\), \(D(z)=z\) for \(R=C\), and \(D(z)=1\) for \(R=A\).
  \item \texttt{Y(d)} --- \texttt{binary potential outcome}; \(\{0,1\}\); \texttt{defines the local complier causal contrast}
  \item \texttt{c\_\allowbreak{}P} --- \texttt{unknown assignment cutoff attached to the law}; \((0,1)\); \texttt{is the discontinuity location}
  \item \texttt{\textbackslash{}\allowbreak{}pi\_\allowbreak{}\{r,\allowbreak{}a,\allowbreak{}b,\allowbreak{}P\}\allowbreak{}} --- \texttt{witness latent response-\allowbreak{}cell regression}; \([0,1]\to\mathbb R\); \texttt{constructs the finite lower-\allowbreak{}bound family and does not define the main minimax class}
  \par\smallskip\noindent\emph{Definition.} For a constructed witness law, the conditional mass at \(X=x\) assigned to \(R=r\), \(Y(0)=a\), and \(Y(1)=b\).
  \item \texttt{P} --- \texttt{full-\allowbreak{}data triangular-\allowbreak{}array law}; \texttt{indexes causal data-\allowbreak{}generating processes through observable reduced forms}
  \par\smallskip\noindent\emph{Definition.} A law of the declared potential variables and observables under threshold assignment at \(c_P\).
  \item \texttt{Z} --- \texttt{threshold assignment}; \(\{0,1\}\); \texttt{is the discontinuous instrument}
  \par\smallskip\noindent\emph{Definition.} \(Z=\mathbf 1\{X\ge c_P\}\).
  \item \texttt{D} --- \texttt{observed treatment receipt}; \(\{0,1\}\); \texttt{is the endogenous treatment}
  \par\smallskip\noindent\emph{Definition.} \(D=D(Z)\).
  \item \texttt{Y} --- \texttt{observed outcome}; \(\{0,1\}\); \texttt{is the observed causal response}
  \par\smallskip\noindent\emph{Definition.} \(Y=Y(D)\).
  \item \texttt{O\_\allowbreak{}i} --- \texttt{observed unit}; \([0,1]\times\{0,1\}^2\); \texttt{is the sampling unit}
  \par\smallskip\noindent\emph{Definition.} \(O_i=(X_i,D_i,Y_i)\), for \(i=1,\ldots,n\).
  \item \texttt{p\_\allowbreak{}\{d,\allowbreak{}y,\allowbreak{}P\}\allowbreak{}} --- \texttt{observable conditional cell regression}; \([0,1]\to[0,1]\); \texttt{certifies positivity inside the explicit witness}
  \par\smallskip\noindent\emph{Definition.} \(p_{d,y,P}(x)=P(D=d,Y=y\mid X=x)\).
  \item \texttt{\textbackslash{}\allowbreak{}mu\_\allowbreak{}\{V,\allowbreak{}P\}\allowbreak{}\textasciicircum{}\{\textbackslash{}\allowbreak{}ell\}\allowbreak{}} --- \texttt{one-\allowbreak{}sided observable reduced-\allowbreak{}form regression}; \(\{x\in[0,1]:x<c_P\}\to[0,1]\) when \(\ell=-\) and \(\{x\in[0,1]:x\ge c_P\}\to[0,1]\) when \(\ell=+\); \texttt{defines the smooth observable primitives used by both upper bounds}
  \par\smallskip\noindent\emph{Definition.} The observable conditional-mean regression restricted to the indicated side of \(c_P\), with \(\mu_{V,P}^{\ell}(c_P)\) denoting the corresponding one-sided limit.
  \item \texttt{d\_\allowbreak{}P} --- \texttt{compliance discontinuity}; \((0,1]\); \texttt{is the local first-\allowbreak{}stage denominator}
  \par\smallskip\noindent\emph{Definition.} \(d_P=\mu_{D,P}^{+}(c_P)-\mu_{D,P}^{-}(c_P)\).
  \item \texttt{y\_\allowbreak{}P} --- \texttt{reduced-\allowbreak{}form outcome discontinuity}; \([-1,1]\); \texttt{is the local intent-\allowbreak{}to-\allowbreak{}treat numerator}
  \par\smallskip\noindent\emph{Definition.} \(y_P=\mu_{Y,P}^{+}(c_P)-\mu_{Y,P}^{-}(c_P)\).
  \item \texttt{\textbackslash{}\allowbreak{}Theta} --- \texttt{effect-\allowbreak{}label space}; \([-1,1]\); \texttt{is the parameter space for local complier effects}
  \par\smallskip\noindent\emph{Definition.} \(\Theta=[-1,1]\).
  \item \texttt{\textbackslash{}\allowbreak{}theta\_\allowbreak{}P} --- \texttt{fuzzy-\allowbreak{}RD local average treatment effect}; \(\Theta\); \texttt{is the causal target}
  \par\smallskip\noindent\emph{Definition.} \(\theta_P=y_P/d_P\).
  \item \texttt{\textbackslash{}\allowbreak{}rho\_\allowbreak{}n} --- \texttt{known-\allowbreak{}location nonparametric noise scale}; \((0,1]\); \texttt{normalizes weak-\allowbreak{}denominator interval length}
  \par\smallskip\noindent\emph{Definition.} \(\rho_n=n^{-s/(2s+1)}\).
  \item \texttt{T\_\allowbreak{}n} --- \texttt{cutoff-\allowbreak{}localization information index}; \((0,\infty)\); \texttt{compares local jump information with location entropy}
  \par\smallskip\noindent\emph{Definition.} \(T_n=n\kappa_n^{2+1/s}\).
  \item \texttt{\textbackslash{}\allowbreak{}mathcal P\_\allowbreak{}n} --- \texttt{triangular class of observable-\allowbreak{}smooth monotone fuzzy-\allowbreak{}RD laws}; \texttt{is the minimax parameter space}
  \item \texttt{I\_\allowbreak{}n} --- \texttt{measurable connected confidence-\allowbreak{}interval rule}; \(([0,1]\times\{0,1\}^2)^n\to\{\text{connected subsets of }\Theta\}\); \texttt{is the decision rule over which minimax length is optimized}
  \item \texttt{I\_\allowbreak{}n\textasciicircum{}\{\textbackslash{}\allowbreak{}mathrm\{known\}\allowbreak{}\}\allowbreak{}} --- \texttt{measurable known-\allowbreak{}cutoff connected confidence-\allowbreak{}interval rule}; \(\bigl(([0,1]\times\{0,1\}^2)^n\times[1/4,3/4]\bigr)\to\{\text{connected subsets of }\Theta\}\); \texttt{is the oracle decision rule that receives the supplied cutoff as an input}
  \item \texttt{L\_\allowbreak{}n} --- \texttt{minimax fixed-\allowbreak{}level expected connected-\allowbreak{}interval length}; \([0,2]\); \texttt{is the headline statistical frontier}
  \item \texttt{L\_\allowbreak{}n\textasciicircum{}\{\textbackslash{}\allowbreak{}mathrm\{known\}\allowbreak{}\}\allowbreak{}} --- \texttt{known-\allowbreak{}cutoff minimax expected connected-\allowbreak{}interval length}; \([0,2]\); \texttt{is the oracle benchmark}
  \item \texttt{a\_\allowbreak{}\textbackslash{}\allowbreak{}star} --- \texttt{effect-\allowbreak{}label half-\allowbreak{}separation}; \((0,1)\); \texttt{fixes the parameter-\allowbreak{}exact high-\allowbreak{}separation pair of causal labels}
  \par\smallskip\noindent\emph{Definition.} \(a_\star=1-v\).
  \item \texttt{b} --- \texttt{common potential-\allowbreak{}outcome discordance mass}; \((0,1)\); \texttt{keeps every parameter-\allowbreak{}exact common potential-\allowbreak{}outcome cell strictly positive}
  \par\smallskip\noindent\emph{Definition.} \(b=1-v/2\).
  \item \texttt{\textbackslash{}\allowbreak{}psi} --- \texttt{compactly supported quartic compliance profile}; \(\mathbb R\to[0,1]\); \texttt{generates smooth translated first-\allowbreak{}stage jumps}
  \item \texttt{g\_\allowbreak{}\textbackslash{}\allowbreak{}star} --- \texttt{smooth common outcome-\allowbreak{}effect profile}; \([0,1]\to[-a_\star,a_\star]\); \texttt{assigns opposite effect labels on two cutoff plateaus}
  \item \texttt{q\_\allowbreak{}\{a,\allowbreak{}b\}\allowbreak{}} --- \texttt{common potential-\allowbreak{}outcome cell table}; \([0,1]\to[0,1]\); \texttt{makes the conditional outcome law invariant across cutoff alternatives}
  \item \texttt{m\_\allowbreak{}\{d,\allowbreak{}y\}\allowbreak{}} --- \texttt{common conditional outcome margin}; \([0,1]\to(0,1)\); \texttt{is the outcome margin shared by all response types and cutoff alternatives}
  \par\smallskip\noindent\emph{Definition.} \(m_{d,y}(x)=\sum_{a_0,a_1\in\{0,1\}}\mathbf 1\{(1-d)a_0+da_1=y\}q_{a_0,a_1}(x)\).
  \item \texttt{R\_\allowbreak{}\textbackslash{}\allowbreak{}star} --- \texttt{outcome-\allowbreak{}invariant comparison law}; \texttt{dominates the translated witness mixtures}
  \item \texttt{h\_\allowbreak{}n} --- \texttt{least-\allowbreak{}favorable bump width}; \((0,\infty)\); \texttt{balances compliance height and smoothness}
  \par\smallskip\noindent\emph{Definition.} \(h_n=(32\kappa_n/M)^{1/s}\).
  \item \texttt{J\_\allowbreak{}n} --- \texttt{number of translated alternatives per effect label}; \(\mathbb N_0\); \texttt{measures cutoff-\allowbreak{}location multiplicity}
  \par\smallskip\noindent\emph{Definition.} \(J_n=\lfloor(16h_n)^{-1}\rfloor\).
  \item \texttt{\textbackslash{}\allowbreak{}mathcal W\_\allowbreak{}n} --- \texttt{finite outcome-\allowbreak{}invariant location-\allowbreak{}mixture witness family}; \texttt{certifies the unknown-\allowbreak{}cutoff lower bound}
  \item \texttt{\textbackslash{}\allowbreak{}mathsf H\_\allowbreak{}\{\textbackslash{}\allowbreak{}mathrm\{crit\}\allowbreak{}\}\allowbreak{}} --- \texttt{critical-\allowbreak{}window derivation handle}; \texttt{specifies the noncommittal route for deriving the exact frontier}
  \item \texttt{I\_\allowbreak{}n\textasciicircum{}\{\textbackslash{}\allowbreak{}mathrm\{SI\}\allowbreak{},\allowbreak{}1\}\allowbreak{}} --- \texttt{Lipschitz four-\allowbreak{}fold scan-\allowbreak{}-\allowbreak{}refine-\allowbreak{}-\allowbreak{}invert connected interval}; is the total explicit interval rule proved to attain the oracle order on the face \(s=1\)
  \item \texttt{\textbackslash{}\allowbreak{}widehat c\_\allowbreak{}\{\textbackslash{}\allowbreak{}mathrm\{prof\}\allowbreak{}\}\allowbreak{}} --- \texttt{profile-\allowbreak{}selected cutoff center}; \([1/4,3/4]\); \texttt{is the cutoff coordinate returned by the higher-\allowbreak{}smoothness common-\allowbreak{}versus-\allowbreak{}piecewise affine profile}
  \item \texttt{\textbackslash{}\allowbreak{}mathsf F\_\allowbreak{}\{\textbackslash{}\allowbreak{}mathrm\{prof\}\allowbreak{}\}\allowbreak{}} --- \texttt{profile-\allowbreak{}construction fallback indicator}; \(\{0,1\}\); \texttt{makes singular-\allowbreak{}design and weak-\allowbreak{}profile fallbacks an explicit measurable output}
  \item \texttt{I\_\allowbreak{}n\textasciicircum{}\{\textbackslash{}\allowbreak{}mathrm\{SI\}\allowbreak{}\}\allowbreak{}} --- \texttt{smoothness-\allowbreak{}indexed total connected confidence-\allowbreak{}interval rule}; dispatches to the proved rectangular Lipschitz rule at \(s=1\) and to the explicit but performance-unresolved profile branch at \(s>1\)
\end{itemize}

\section{Assumptions}

\begin{assumption}[ass:iid-sampling]\label{ass:iid-sampling}
The observations \(O_1,\ldots,O_n\) are independent and identically distributed under \(P\).
\par\smallskip\noindent\emph{Standard} (Independent random-design sampling, \cite{NoackRothe2024}).
\end{assumption}

\begin{assumption}[ass:uniform-running-variable]\label{ass:uniform-running-variable}
The marginal law of \(X\) under \(P\) is uniform on \([0,1]\).
\par\smallskip\noindent\emph{Standard} (Uniform random-design specialization, \cite{PorterYu2015}).
\end{assumption}

\begin{assumption}[ass:interior-cutoff]\label{ass:interior-cutoff}
The cutoff satisfies \(c_P\in[1/4,3/4]\).
\par\smallskip\noindent\emph{Standard} (Interior cutoff trimming, \cite{PorterYu2015}).
\end{assumption}

\begin{assumption}[ass:potential-reduced-form-continuity]\label{ass:potential-reduced-form-continuity}
For each \(z\in\{0,1\}\), the maps \(x\mapsto E_P[D(z)\mid X=x]\) and \(x\mapsto E_P[Y(D(z))\mid X=x]\) are continuous at \(c_P\).
\par\smallskip\noindent\emph{Standard} (Fuzzy-RD continuity of potential reduced forms, \cite{HahnToddVanDerKlaauw2001}).
\end{assumption}

\begin{assumption}[ass:observable-holder-extensions]\label{ass:observable-holder-extensions}
For every \(V\in\{D,Y\}\) and \(\ell\in\{-,+\}\), the one-sided regression \(\mu_{V,P}^{\ell}\) admits an extension to \([0,1]\) in the radius-\(M\) H\"older ball of order \(s\), with the derivative interpreted in the usual way when \(s>1\).
\par\smallskip\noindent\emph{Standard} (Adapted/sufficient specialization, cf. NoackRothe2024 (one-sided observable H{\"o}lder extensions), \cite{NoackRothe2024}).
\end{assumption}

\begin{assumption}[ass:local-compliance-strength]\label{ass:local-compliance-strength}
The compliance discontinuity satisfies \(\kappa_n\le d_P\le2\kappa_n\).
\par\smallskip\noindent\emph{Novel.} The factor-two shell is a deliberate strength-indexed triangular slice that defines uniform minimax risk at the scale \(\kappa_n\); it is not attributed to a published fuzzy-RD assumption. It allows arbitrary local-to-zero strength sequences and groups laws whose compliance jumps have the same empirical order, as occurs when take-up contrasts weaken across increasingly local study populations.
\end{assumption}

\begin{assumption}[ass:observable-outcome-variance-floor]\label{ass:observable-outcome-variance-floor}
For each \(\ell\in\{-,+\}\) and every \(x\) on side \(\ell\) of \(c_P\), \(\operatorname{Var}_P(Y\mid X=x,\ell)\ge v\).
\par\smallskip\noindent\emph{Standard} (Adapted/sufficient specialization, cf. NoackRothe2024 (binary-outcome variance floor), \cite{NoackRothe2024}).
\end{assumption}

\begin{assumption}[ass:observable-correlation-separation]\label{ass:observable-correlation-separation}
For each \(\ell\in\{-,+\}\) and every \(x\) on side \(\ell\) of \(c_P\), \[|\operatorname{Cov}_P(Y,D\mid X=x,\ell)|\le(1-v)\sqrt{\operatorname{Var}_P(Y\mid X=x,\ell)\operatorname{Var}_P(D\mid X=x,\ell)}.\] If \(\operatorname{Var}_P(D\mid X=x,\ell)=0\), then the covariance is zero and the display is interpreted as \(0\le0\), without forming a correlation coefficient.
\par\smallskip\noindent\emph{Standard} (Adapted/sufficient specialization, cf. NoackRothe2024 (observable correlation separation), \cite{NoackRothe2024}).
\end{assumption}

\section{Definitions}

\begin{definition}[def:model-class]\label{def:model-class}
\par\smallskip\noindent\emph{Defined object.} \texttt{\textbackslash{}\allowbreak{}mathcal P\_\allowbreak{}n}
\par\smallskip\noindent\emph{Construction.} \[\mathcal P_n(s,M,v,\kappa_n)=\{P:\text{ properties ass:iid-sampling, ass:uniform-running-variable, ass:interior-cutoff, ass:potential-reduced-form-continuity, ass:observable-holder-extensions, ass:local-compliance-strength, ass:observable-outcome-variance-floor, and ass:observable-correlation-separation hold at index }n\}.\]
\par\smallskip\noindent\emph{Member properties.} \texttt{ass:\allowbreak{}iid-\allowbreak{}sampling}; \texttt{ass:\allowbreak{}uniform-\allowbreak{}running-\allowbreak{}variable}; \texttt{ass:\allowbreak{}interior-\allowbreak{}cutoff}; \texttt{ass:\allowbreak{}potential-\allowbreak{}reduced-\allowbreak{}form-\allowbreak{}continuity}; \texttt{ass:\allowbreak{}observable-\allowbreak{}holder-\allowbreak{}extensions}; \texttt{ass:\allowbreak{}local-\allowbreak{}compliance-\allowbreak{}strength}; \texttt{ass:\allowbreak{}observable-\allowbreak{}outcome-\allowbreak{}variance-\allowbreak{}floor}; \texttt{ass:\allowbreak{}observable-\allowbreak{}correlation-\allowbreak{}separation}.
\end{definition}

\begin{definition}[def:minimax-length]\label{def:minimax-length}
\par\smallskip\noindent\emph{Defined object.} \texttt{L\_\allowbreak{}n}
\par\smallskip\noindent\emph{Construction.} \[L_n(\kappa_n)=\inf_{I_n}\sup_{P\in\mathcal P_n}E_P[|I_n(O_1,\ldots,O_n)|],\qquad \inf_{P\in\mathcal P_n}P\{\theta_P\in I_n(O_1,\ldots,O_n)\}\ge1-\alpha,\] where the infimum ranges over measurable connected intervals contained in \(\Theta\).
\par\smallskip\noindent\emph{Inputs.} \texttt{def:\allowbreak{}model-\allowbreak{}class}.
\end{definition}

\begin{definition}[def:known-cutoff-length]\label{def:known-cutoff-length}
\par\smallskip\noindent\emph{Defined object.} \texttt{L\_\allowbreak{}n\textasciicircum{}\{\textbackslash{}\allowbreak{}mathrm\{known\}\allowbreak{}\}\allowbreak{}}
\par\smallskip\noindent\emph{Construction.} \[L_n^{\mathrm{known}}(\kappa_n)=\inf_{I_n^{\mathrm{known}}}\sup_{P\in\mathcal P_n}E_P[|I_n^{\mathrm{known}}((O_1,\ldots,O_n),c_P)|],\qquad \inf_{P\in\mathcal P_n}P\{\theta_P\in I_n^{\mathrm{known}}((O_1,\ldots,O_n),c_P)\}\ge1-\alpha,\] where the infimum ranges only over measurable rules of the declared known-cutoff oracle type with connected values contained in \(\Theta\).
\par\smallskip\noindent\emph{Inputs.} \texttt{def:\allowbreak{}model-\allowbreak{}class}.
\end{definition}

\begin{definition}[def:quartic-profile]\label{def:quartic-profile}
\par\smallskip\noindent\emph{Defined object.} \texttt{\textbackslash{}\allowbreak{}psi}
\par\smallskip\noindent\emph{Construction.} \[\psi(u)=(1-u^2)^2\mathbf 1\{|u|\le1\}.\]
\end{definition}

\begin{definition}[def:common-outcome-profile]\label{def:common-outcome-profile}
\par\smallskip\noindent\emph{Defined object.} \texttt{g\_\allowbreak{}\textbackslash{}\allowbreak{}star}
\par\smallskip\noindent\emph{Construction.} Let \(B(t)=1-10t^3+15t^4-6t^5\). Define \(g_\star(x)=a_\star\) on \([0,3/8]\), \(g_\star(x)=a_\star\{2B(4x-3/2)-1\}\) on \((3/8,5/8)\), and \(g_\star(x)=-a_\star\) on \([5/8,1]\).
\end{definition}

\begin{definition}[def:common-outcome-table]\label{def:common-outcome-table}
\par\smallskip\noindent\emph{Defined object.} \texttt{q\_\allowbreak{}\{a,\allowbreak{}b\}\allowbreak{}}
\par\smallskip\noindent\emph{Construction.} Set \(b=1-v/2\).  For \(x\in[0,1]\), define \(q_{00}(x)=q_{11}(x)=(1-b)/2\), \(q_{01}(x)=\{b+g_\star(x)\}/2\), and \(q_{10}(x)=\{b-g_\star(x)\}/2\).
\par\smallskip\noindent\emph{Inputs.} \texttt{def:\allowbreak{}common-\allowbreak{}outcome-\allowbreak{}profile}.
\end{definition}

\begin{definition}[def:outcome-invariant-reference]\label{def:outcome-invariant-reference}
\par\smallskip\noindent\emph{Defined object.} \texttt{R\_\allowbreak{}\textbackslash{}\allowbreak{}star}
\par\smallskip\noindent\emph{Construction.} Under \(R_\star\), \(X\sim\operatorname{Unif}[0,1]\), \(R_\star(D=d\mid X=x)=1/2\), and \(R_\star(Y=y\mid D=d,X=x)=m_{d,y}(x)\). This comparison law is used only as a dominating measure for the two in-class mixtures.
\par\smallskip\noindent\emph{Inputs.} \texttt{def:\allowbreak{}common-\allowbreak{}outcome-\allowbreak{}table}.
\end{definition}

\begin{definition}[def:outcome-invariant-family]\label{def:outcome-invariant-family}
\par\smallskip\noindent\emph{Defined object.} \texttt{\textbackslash{}\allowbreak{}mathcal W\_\allowbreak{}n}
\par\smallskip\noindent\emph{Construction.} For \(j=1,\ldots,J_n\), set \(c_{j,+}=1/4+(2j-1)h_n\) and \(c_{j,-}=5/8+(2j-1)h_n\). For \(\ell\in\{+,-\}\), put \(w_{j,\ell}(x)=\kappa_n\psi((x-c_{j,\ell})/h_n)\), \(\pi_{C,a,b,P_{j,\ell}}(x)=w_{j,\ell}(x)q_{a,b}(x)\), and \(\pi_{A,a,b,P_{j,\ell}}(x)=\pi_{N,a,b,P_{j,\ell}}(x)=\{1-w_{j,\ell}(x)\}q_{a,b}(x)/2\). The induced laws form \(\mathcal W_n=\{P_{j,\ell}:1\le j\le J_n,\ell\in\{+,-\}\}\), each with cutoff \(c_{j,\ell}\).
\par\smallskip\noindent\emph{Inputs.} \texttt{def:\allowbreak{}quartic-\allowbreak{}profile}; \texttt{def:\allowbreak{}common-\allowbreak{}outcome-\allowbreak{}table}.
\end{definition}

\begin{definition}[def:scan-invert]\label{def:scan-invert}
\par\smallskip\noindent\emph{Defined object.} \texttt{I\_\allowbreak{}n\textasciicircum{}\{\textbackslash{}\allowbreak{}mathrm\{SI\}\allowbreak{}\}\allowbreak{}}
\par\smallskip\noindent\emph{Construction.} If \(s=1\), define \(I_n^{\mathrm{SI}}=I_n^{\mathrm{SI},1}\) by
\(\mathrm{def:lipschitz\mbox{-}scan\mbox{-}invert}\).  Suppose instead
that \(1<s\le2\), and set \(m=\lfloor n/4\rfloor\).  If \(m=0\),
return \(I_n^{\mathrm{SI}}=\Theta\) before forming folds or a
\(1/m\) statistic.  Otherwise obtain
\((\widehat c_{\mathrm{prof}},\mathsf F_{\mathrm{prof}})\) from
\(\mathrm{def:profile\mbox{-}cutoff\mbox{-}center}\).  If
\(\mathsf F_{\mathrm{prof}}=1\), return \(\Theta\).  On the remaining
branch put
\[
 K_1(u)=(4-6u)\mathbf 1\{0\le u\le1\},\qquad
 b_s=m^{-1/(2s+1)},\qquad L_\alpha=\log(32/\alpha),
\]
\[
 R_{\mathrm{prof}}
 =\frac{2^{16}L_\alpha}{m\kappa_n^2}+\frac1n,
 \qquad
 r_{\mathrm{prof}}
 =8Mb_s^s+16MR_{\mathrm{prof}}
  +8\|K_1\|_\infty
     \sqrt{\frac{L_\alpha}{mb_s}}.
\]
If \(b_s>1/8\), return \(\Theta\) before forming a fourth-fold
contrast.  On the remaining branch, using the fourth fold and the
fixed denominator \(mb_s\), define
\[
 \widehat\Delta_{4,b_s}^{V}(t)
 =\frac1{mb_s}\sum_{i\in F_4}V_i
 \left\{K_1\!\left(\frac{X_i-t}{b_s}\right)
       -K_1\!\left(\frac{t-X_i}{b_s}\right)\right\},
 \quad V\in\{D,Y\},
\]
and set \(\widehat y=\widehat\Delta_{4,b_s}^{Y}
(\widehat c_{\mathrm{prof}})\) and
\(\widehat d=\widehat\Delta_{4,b_s}^{D}
(\widehat c_{\mathrm{prof}})\).  If
\(R_{\mathrm{prof}}>b_s/8\), \(r_{\mathrm{prof}}>\kappa_n/8\), or
\(\widehat d<\kappa_n/2\), return \(\Theta\).  Otherwise return
\[
 I_n^{\mathrm{SI}}
 =\operatorname{hull}\left\{\theta\in\Theta:
 |\widehat y-\theta\widehat d|\le2r_{\mathrm{prof}}\right\}.
\]
The higher-smoothness branch is a completely specified measurable
sample map; no performance assertion for that branch is part of this
definition.
\par\smallskip\noindent\emph{Inputs.} \texttt{def:\allowbreak{}model-\allowbreak{}class}; \texttt{def:\allowbreak{}lipschitz-\allowbreak{}scan-\allowbreak{}invert}; \texttt{def:\allowbreak{}profile-\allowbreak{}cutoff-\allowbreak{}center}.
\end{definition}

\begin{definition}[def:critical-handle]\label{def:critical-handle}
\par\smallskip\noindent\emph{Defined object.} \texttt{\textbackslash{}\allowbreak{}mathsf H\_\allowbreak{}\{\textbackslash{}\allowbreak{}mathrm\{crit\}\allowbreak{}\}\allowbreak{}}
\par\smallskip\noindent\emph{Construction.} \texttt{Derive bias-\allowbreak{}compensated local likelihoods indexed jointly by cutoff and effect label;\allowbreak{} solve the normalized optimal-\allowbreak{}recovery problem for the compliance and outcome compensation profiles;\allowbreak{} truncate and renormalize the translated-\allowbreak{}cutoff mixtures at the support boundary;\allowbreak{} retain the effect-\allowbreak{}label coordinate in the testing experiment;\allowbreak{} and invert the resulting tests before taking the shortest honest connected hull.}
\par\smallskip\noindent\emph{Inputs.} \texttt{def:\allowbreak{}minimax-\allowbreak{}length}; \texttt{def:\allowbreak{}outcome-\allowbreak{}invariant-\allowbreak{}family}.
\end{definition}

\begin{definition}[def:lipschitz-scan-invert]\label{def:lipschitz-scan-invert}
\par\smallskip\noindent\emph{Defined object.} \texttt{I\_\allowbreak{}n\textasciicircum{}\{\textbackslash{}\allowbreak{}mathrm\{SI\}\allowbreak{},\allowbreak{}1\}\allowbreak{}}
\par\smallskip\noindent\emph{Construction.} Set \(m=\lfloor n/4\rfloor\).  If \(m=0\), define
\(I_n^{\mathrm{SI},1}(O_1,\ldots,O_n)=\Theta\) and stop; in
particular, this branch is taken before forming folds, before using a
\(1/m\) statistic, and before defining \(b=m^{-1/3}\).  If \(m\ge1\),
form the four deterministic folds
\(F_j=\{(j-1)m+1,\ldots,jm\}\), discarding observations with index
larger than \(4m\), and put
\[
 K_0(u)=\mathbf 1\{0\le u\le1\},\qquad
 H=\frac{\kappa_n}{2^{12}M}.
\]
For \(V\in\{D,Y\}\), \(j\in\{1,2,3,4\}\), \(h>0\), and any
\(t\) for which \([t-h,t+h]\subset[0,1]\), define the fixed-denominator
contrast
\[
 \widehat\Delta_{j,h}^{V}(t)
 =\frac1{mh}\sum_{i\in F_j}V_i
 \left\{K_0\!\left(\frac{X_i-t}{h}\right)
       -K_0\!\left(\frac{t-X_i}{h}\right)\right\}.
\]
Let
\[
 \mathcal G_H=\left\{\frac14+\frac{kH}{64}:
 0\le k\le\left\lfloor\frac{32}{H}\right\rfloor\right\},
 \qquad
 \mathcal S_H=\left\{t\in\mathcal G_H:
 \widehat\Delta_{1,H}^{D}(t)
 \ge\max_{q\in\mathcal G_H}\widehat\Delta_{1,H}^{D}(q)
       -\frac{\kappa_n}{8}\right\},
\]
and let \(B_H\) be the union, intersected with \([1/4,3/4]\), of the
closed \(H\)-neighborhoods of the points in \(\mathcal S_H\).  On each
closed component of \(B_H\), take the deterministic grid that includes
both endpoints and has mesh at most \(1/n\); let \(\mathcal G_n(B_H)\)
be the union of these grids.  Define \(\widehat c\) as the smallest
maximizer over \(\mathcal G_n(B_H)\) of
\[
 t\longmapsto
 \overline\Delta_H^D(t)
 =\frac12\left\{\widehat\Delta_{2,H}^{D}(t)
                 +\widehat\Delta_{3,H}^{D}(t)\right\}.
\]
Now put
\[
 b=m^{-1/3},\qquad L_\alpha=\log(32/\alpha),\qquad
 R_c=\frac{2^{16}L_\alpha}{m\kappa_n^2}+\frac1n,
\]
\[
 r_n=8Mb+\frac{16\kappa_nR_c}{b}+16MR_c
     +8\|K_0\|_\infty\sqrt{\frac{L_\alpha}{mb}},
\]
If \(b>1/8\), return \(I_n^{\mathrm{SI},1}=\Theta\) before forming a
fold-four contrast.  On the remaining branch set
\[
 \widehat y=\widehat\Delta_{4,b}^{Y}(\widehat c),
 \qquad
 \widehat d=\widehat\Delta_{4,b}^{D}(\widehat c).
\]
If \(R_c>H/8\) or \(r_n>\kappa_n/8\), return
\(I_n^{\mathrm{SI},1}=\Theta\).  If those checks pass but
\(\widehat d<\kappa_n/2\), again return \(\Theta\).  Otherwise return
\[
 I_n^{\mathrm{SI},1}
 =\operatorname{hull}\left\{\theta\in\Theta:
       |\widehat y-\theta\widehat d|\le2r_n\right\},
\]
with the hull of the empty set defined to be empty.  All maximizers use
the displayed smallest-argument tie breaking, so this is a measurable
map of the sample.
\par\smallskip\noindent\emph{Inputs.} \texttt{def:\allowbreak{}model-\allowbreak{}class}.
\end{definition}

\begin{definition}[def:profile-cutoff-center]\label{def:profile-cutoff-center}
\par\smallskip\noindent\emph{Defined object.} \texttt{(\textbackslash{}\allowbreak{}widehat c\_\allowbreak{}\{\textbackslash{}\allowbreak{}mathrm\{prof\}\allowbreak{}\}\allowbreak{},\allowbreak{}\textbackslash{}\allowbreak{}mathsf F\_\allowbreak{}\{\textbackslash{}\allowbreak{}mathrm\{prof\}\allowbreak{}\}\allowbreak{})}
\par\smallskip\noindent\emph{Construction.} Set \(m=\lfloor n/4\rfloor\).  If \(m=0\), return
\((\widehat c_{\mathrm{prof}},\mathsf F_{\mathrm{prof}})=(1/2,1)\)
and stop; this branch precedes the definition of \(b=m^{-1/3}\), every
\(1/m\) statistic, and every fold construction.  If \(m\ge1\), put
\[
 b=m^{-1/3},\qquad
 H=\left(\frac{\kappa_n}{2^{12}M}\right)^{1/s},\qquad
 \delta=\frac{\kappa_n}{32},
\]
form \(F_j=\{(j-1)m+1,\ldots,jm\}\), \(j=1,\ldots,4\), and take a
deterministic grid \(\mathcal T_n\subset[1/4,3/4]\) containing both
endpoints and having mesh at most
\(\eta_n=\min\{n^{-1},bH/64\}\).  For \(\ell\in\{-,+\}\), define
\[
 z_{t,-}(x)=\left(1,\frac{t-x}{H}\right)^\top,
 \quad A_{t,-}=[t-H,t),\qquad
 z_{t,+}(x)=\left(1,\frac{x-t}{H}\right)^\top,
 \quad A_{t,+}=[t,t+H],
\]
Here \(H\le2^{-8}<1/4\), by \(\kappa_n\le1/8\), \(M\ge2\), and
\(1<s\le2\), so every displayed flank is contained in \([0,1]\).
For \(j\in\{1,2\}\), put
\[
 \widehat G_{j,t}^{\ell}
 =\frac1m\sum_{i\in F_j}
   \mathbf 1\{X_i\in A_{t,\ell}\}
   z_{t,\ell}(X_i)z_{t,\ell}(X_i)^\top.
\]
Under the uniform design the corresponding population matrix is
\[
 G_H=H\begin{pmatrix}1&1/2\\[2pt]1/2&1/3\end{pmatrix},
 \qquad
 \lambda_{\min}(G_H)=\frac{H(4-\sqrt{13})}{6}.
\]
If any displayed empirical flank matrix has smallest eigenvalue below
\(H(4-\sqrt{13})/12\), return \((1/2,1)\).

Otherwise, on each training fold \(j\in\{1,2\}\) and at every
\(t\in\mathcal T_n\), compute the lexicographically smallest constrained
least-squares fits of \(D\): a common affine fit on
\([t-H,t+H]\) and two separate affine fits on \(A_{t,-}\) and
\(A_{t,+}\).  Each affine fit has intercept in \([-b,1+b]\) and slope
in \([-M-b,M+b]\); these compact constraints make every fit total and
measurable.  Denote the common and piecewise fitted functions by
\(\widehat g_{j,t}^{0}\) and \(\widehat g_{j,t}^{1}\), respectively.
For \((j,k)=(1,2)\) and \((2,3)\), define the held-out improvement
\[
 \widehat Q_{j\to k}(t)=\frac1m\sum_{i\in F_k}
 \mathbf 1\{|X_i-t|\le H\}
 \left[
   \{D_i-\widehat g_{j,t}^{0}(X_i)\}^2
  -\{D_i-\widehat g_{j,t}^{1}(X_i)\}^2
 \right],
 \qquad
 \widehat Q(t)=\frac{\widehat Q_{1\to2}(t)+
                           \widehat Q_{2\to3}(t)}2.
\]
Let \(\widetilde c\) be the smallest maximizer of \(\widehat Q\).  Let
\(\widehat d_{\mathrm{prof}}(\widetilde c)\) be the average, over the
two training folds, of the fitted right intercept minus fitted left
intercept at \(\widetilde c\).  If
\(\widehat d_{\mathrm{prof}}(\widetilde c)<\kappa_n-8\delta\), return
\((\widetilde c,1)\); otherwise return \((\widetilde c,0)\).  Thus the
fallback indicator \(\mathsf F_{\mathrm{prof}}\) is an explicit
component of the returned object.
\par\smallskip\noindent\emph{Inputs.} \texttt{def:\allowbreak{}model-\allowbreak{}class}.
\end{definition}

\section{Main results}

\begin{proposition}[prop:wald-late-identification]\label{prop:wald-late-identification}
For every \(P\in\mathcal P_n\), continuity of the potential reduced forms, monotonicity encoded by \(R\in\{N,C,A\}\), consistency, and exclusion imply \[d_P=P(R=C\mid X=c_P),\qquad y_P=E_P[(Y(1)-Y(0))\mathbf 1\{R=C\}\mid X=c_P].\] Therefore \(\theta_P=y_P/d_P=E_P[Y(1)-Y(0)\mid R=C,X=c_P]\), where all conditional quantities at \(c_P\) are their continuous versions.
\end{proposition}
\begin{proof}
Fix \(P\in\mathcal P_n\).  By the definitions of threshold assignment and observed treatment, the right and left limits of the treatment regression are, respectively,
\[
 \mu^+_{D,P}(c_P)=E_P[D(1)\mid X=c_P],\qquad
 \mu^-_{D,P}(c_P)=E_P[D(0)\mid X=c_P].
\]
Here the equalities at \(c_P\) use the continuous versions supplied by assumption ass:potential-reduced-form-continuity.  The response-type definition gives the pointwise identity
\[
 D(1)-D(0)=\mathbf 1\{R=C\}.
\]
Subtracting the two preceding conditional expectations therefore yields
\[
 d_P=E_P[D(1)-D(0)\mid X=c_P]=P(R=C\mid X=c_P).
\]
Likewise, consistency and exclusion are encoded by \(Y=Y(D)\): immediately to the right of the cutoff the observed potential response is \(Y(D(1))\), and immediately to the left it is \(Y(D(0))\).  Hence continuity of the potential reduced forms gives
\[
 y_P=E_P[Y(D(1))-Y(D(0))\mid X=c_P].
\]
For never-takers and always-takers, \(D(1)=D(0)\), whereas for compliers \((D(0),D(1))=(0,1)\).  Consequently, pointwise,
\[
 Y(D(1))-Y(D(0))=(Y(1)-Y(0))\mathbf 1\{R=C\},
\]
and thus
\[
 y_P=E_P[(Y(1)-Y(0))\mathbf 1\{R=C\}\mid X=c_P].
\]
Assumption ass:local-compliance-strength gives \(d_P\ge\kappa_n>0\), so division is valid.  Dividing the last display by the formula for \(d_P\) proves
\[
 \theta_P=\frac{y_P}{d_P}=E_P[Y(1)-Y(0)\mid R=C,X=c_P].
\]
\end{proof}
\par\smallskip\noindent \textit{Justification.} This proposition connects the observable jump ratio that drives inference to the complier causal effect without imposing latent-cell smoothness on the minimax class. \textit{Gap.} HahnToddVanDerKlaauw2001 supplies the standard fuzzy-RD Wald interpretation; the proposition records that interpretation for the observable triangular class used by the frontier. \textit{Consumer.} Every interval-length theorem can be read as inference for a local complier effect rather than for an abstract ratio of discontinuities.

\begin{theorem}[thm:known-cutoff-benchmark]\label{thm:known-cutoff-benchmark}
There exist constants \(0<c_0<C_0<\infty\), depending only on \(s,M,v,\alpha\), such that for all sufficiently large \(n\), \[c_0\min\{1,\rho_n/\kappa_n\}\le L_n^{\mathrm{known}}(\kappa_n)\le C_0\min\{1,\rho_n/\kappa_n\}.\]
\end{theorem}
\begin{proof}
Put \(c=c_P\), which is supplied to the oracle rule.  We first construct an explicit uniformly honest rule.  Let
\[
 K(u)=(4-6u)\mathbf 1\{0\le u\le1\}.
\]
Then \(\int_0^1K(u)\,du=1\), \(\int_0^1uK(u)\,du=0\), and \(B_K=\int_0^1K(u)^2\,du<\infty\).  Set \(h=n^{-1/(2s+1)}\).  For \(V\in\{D,Y\}\), define the sample maps
\[
 \widehat\mu_V^+=\frac1{nh}\sum_{i=1}^nK\left(\frac{X_i-c}{h}\right)
 \mathbf 1\{c\le X_i\le c+h\}V_i,
\]
\[
 \widehat\mu_V^-=\frac1{nh}\sum_{i=1}^nK\left(\frac{c-X_i}{h}\right)
 \mathbf 1\{c-h\le X_i<c\}V_i.
\]
For all sufficiently large \(n\), assumption ass:interior-cutoff ensures that both windows lie in \([0,1]\).  Uniformity of \(X\) and a change of variables show that their expectations are the corresponding integrals of \(K(u)\mu^\ell_{V,P}(c\pm hu)\).

For \(s=1\), the H\"older condition directly bounds the bias by
\[
 M h\int_0^1|K(u)|u\,du.
\]
For \(1<s\le2\), write \(\beta=s-1\).  The fundamental theorem of calculus and the derivative H\"older bound give
\[
 |f(c+t)-f(c)-tf'(c)|
 \le\int_0^{|t|}M r^\beta\,dr
 =\frac{M}{s}|t|^s.
\]
The zero first moment of \(K\) cancels the linear term, giving the same order \(C_KMh^s\).  Since \(0\le V_i\le1\), independence and the uniform design give
\[
 \operatorname{Var}_P(\widehat\mu_V^\ell)\le\frac{B_K}{nh}.
\]
For any centered random variable \(U\), the pointwise inequality \(\mathbf 1\{|U|>t\}\le U^2/t^2\) implies \(P(|U|>t)\le E[U^2]/t^2\).  Applying this inequality to the four estimators and taking a union bound, there is a constant \(C=C(s,M,\alpha)\) such that, with probability at least \(1-\alpha\), all four errors are at most \(C\rho_n/2\).  Thus, for
\[
 \widehat d=\widehat\mu_D^+-\widehat\mu_D^-,\qquad
 \widehat y=\widehat\mu_Y^+-\widehat\mu_Y^-,\qquad r_n=C\rho_n,
\]
the event
\[
 \mathcal E_P=\{|\widehat d-d_P|\le r_n,\ |\widehat y-y_P|\le r_n\}
\]
has probability at least \(1-\alpha\), uniformly over \(P\in\mathcal P_n\).

If \(\kappa_n\le2r_n\), define the oracle interval to be \(\Theta\).  Otherwise put
\[
 \widetilde d=\max\{\widehat d,\kappa_n/2\},\qquad
 \widetilde\theta=\Pi_\Theta(\widehat y/\widetilde d),
\]
where \(\Pi_\Theta\) is metric projection onto \([-1,1]\), and define
\[
 I_n^{\mathrm{known}}=
 \Theta\cap[\widetilde\theta-4r_n/\kappa_n,
              \widetilde\theta+4r_n/\kappa_n].
\]
This is a measurable connected interval.  On \(\mathcal E_P\) in the second regime, \(\widehat d\ge d_P-r_n>\kappa_n/2\), so \(\widetilde d=\widehat d\).  Since \(|y_P|=|\theta_P|d_P\le d_P\),
\[
 \left|\frac{\widehat y}{\widehat d}-\frac{y_P}{d_P}\right|
 \le\frac{|\widehat y-y_P|}{\widehat d}
 +\frac{|y_P|\,|\widehat d-d_P|}{d_P\widehat d}
 \le\frac{2r_n}{\widehat d}\le\frac{4r_n}{\kappa_n}.
\]
Projection cannot increase distance from \(\theta_P\in\Theta\), so the interval covers \(\theta_P\) on \(\mathcal E_P\).  It is therefore uniformly \((1-\alpha)\)-honest.  Its length is deterministically at most \(8r_n/\kappa_n\) in the second regime and is two in the first.  Because \(r_n=C\rho_n\), enlarging a constant if necessary gives
\[
 L_n^{\mathrm{known}}(\kappa_n)\le C_0\min\{1,\rho_n/\kappa_n\}.
\]

For the converse, it is enough to exhibit two laws in the same class with the supplied cutoff \(c=1/2\).  Let
\[
 \varphi(u)=(1-u^2)^2\mathbf 1\{|u|\le1\},\quad
 a_n=\eta\min\{\kappa_n,\rho_n\},\quad
 \delta_n=a_n/\kappa_n,\quad
 h_a=(8a_n/M)^{1/s},
\]
where \(\eta\in(0,1/8]\) will be fixed below.  Under both laws, take \(X\) uniform and, conditionally on \(X\), assign response-type probabilities
\[
 P(R=C)=\kappa_n,\qquad P(R=N)=P(R=A)=\frac{1-\kappa_n}{2}.
\]
For never-takers and always-takers let the four potential-outcome cells each have probability \(1/4\).  Under \(P_t\), \(t\in\{0,1\}\), give compliers the table
\[
 q_{00,t}=q_{11,t}=\frac38,\qquad
 q_{01,t}(x)=\frac{1/4+g_t(x)}2,\qquad
 q_{10,t}(x)=\frac{1/4-g_t(x)}2,\qquad
 g_t(x)=t\delta_n\varphi\left(\frac{x-c}{h_a}\right).
\]
For large \(n\), \(h_a<1/4\); all cells are nonnegative because \(\delta_n\le\eta\le1/8\).  Both laws have
\[
 \mu_D^-=(1-\kappa_n)/2,\qquad \mu_D^+=(1+\kappa_n)/2,
\]
and their outcome extensions are
\[
 \mu_{Y,P_t}^-(x)=\frac12-\frac{\kappa_ng_t(x)}2,\qquad
 \mu_{Y,P_t}^+(x)=\frac12+\frac{\kappa_ng_t(x)}2.
\]
The derivative calculations \(\lVert\varphi'\rVert_\infty\le2\) and \([\varphi']_{s-1}\le8\), together with \(h_a^s=8a_n/M\), show that the nonconstant outcome parts have Lipschitz constant at most \(M/8\) for \(s=1\), and derivative supremum at most \(M/8\) and derivative H\"older seminorm at most \(M/2\) for \(1<s\le2\).  Thus the observable extensions lie in the radius-\(M\) H\"older ball.  The corresponding potential reduced forms are the same continuous functions, so the causal continuity condition holds.  At the cutoff,
\[
 d_{P_t}=\kappa_n,\qquad \theta_{P_0}=0,\qquad \theta_{P_1}=\delta_n.
\]

For completeness, the binary outcome has mean \(1/2\mp\kappa_ng_t/2\) on the two sides, so
\[
 \operatorname{Var}_{P_t}(Y\mid X=x)=\frac14-\frac{\kappa_n^2g_t(x)^2}{4}>\frac{63}{256}>v.
\]
Also \(\operatorname{Var}(D\mid X=x)=(1-\kappa_n^2)/4\ge63/256\), and a direct two-point covariance calculation on either side gives
\[
 |\operatorname{Cov}(Y,D\mid X=x)|=\frac{\kappa_n|g_t(x)|(1-\kappa_n)}4\le\frac1{256}.
\]
The resulting absolute correlation is at most \(1/63<63/64<1-v\).  Thus \(P_0,P_1\in\mathcal P_n\).

Under \(P_0\), each outcome cell within a treatment stratum is equiprobable.  On the left, passing to \(P_1\) changes only the two \(D=0\) cells by \(\pm\kappa_ng_1(x)/2\); on the right it changes only the two \(D=1\) cells by those amounts.  Since the relevant baseline cell probability is \((1+\kappa_n)/4\), the conditional chi-square divergence is at most \(2\kappa_n^2g_1(x)^2\).  Therefore, with \(C_\varphi=\int_{-1}^1\varphi(u)^2\,du=256/315\),
\[
 \chi^2(P_1,P_0)\le2C_\varphi a_n^2h_a.
\]
Product expansion and \(1+x\le e^x\) yield
\[
 \chi^2(P_1^{\otimes n},P_0^{\otimes n})
 \le\exp\left\{2C_\varphi n a_n^2(8a_n/M)^{1/s}\right\}-1.
\]
Since \(n\rho_n^{2+1/s}=1\), the exponent is at most a constant times \(\eta^{2+1/s}\).  Choose \(\eta=\eta(s,M,\alpha)>0\) small enough that the last display is at most \((1-2\alpha)^2\).  Writing \(L=dP_1^{\otimes n}/dP_0^{\otimes n}\), the total-variation identity and Cauchy--Schwarz give
\[
 \operatorname{TV}(P_1^{\otimes n},P_0^{\otimes n})
 =\frac12E_{P_0^{\otimes n}}|L-1|
 \le\frac12\sqrt{\chi^2(P_1^{\otimes n},P_0^{\otimes n})}
 \le\frac{1-2\alpha}{2}.
\]
If an oracle interval is honest under both laws, then under \(P_0\) the probability that it contains both \(0\) and \(\delta_n\) is at least
\[
 1-2\alpha-\operatorname{TV}(P_0^{\otimes n},P_1^{\otimes n})
 \ge\frac{1-2\alpha}{2}.
\]
Connectedness forces its length to be at least \(\delta_n\) on that event.  Hence
\[
 \sup_{P\in\mathcal P_n}E_P[|I_n^{\mathrm{known}}|]
 \ge\frac{\eta(1-2\alpha)}2\min\{1,\rho_n/\kappa_n\}.
\]
Taking the infimum over honest oracle rules proves the lower bound and completes the theorem.
\end{proof}
\par\smallskip\noindent \textit{Justification.} The oracle rate isolates weak-denominator uncertainty from the additional cost of locating the cutoff and is uniform over the observable reduced-form class. \textit{Gap.} NoackRothe2024 proves honest known-cutoff Anderson--Rubin coverage and regular strong-first-stage equivalence, but does not state this triangular-array minimax expected-length order. \textit{Consumer.} The all-regime frontier uses this theorem as its required oracle reduction when the cutoff is supplied or strongly localized.

\begin{openendedquestion}[oeq:critical-window-frontier]\label{oeq:critical-window-frontier}
Using \(\mathsf H_{\mathrm{crit}}\), determine a nonvariational
fixed-\(\alpha\) expression for \(L_n(\kappa_n)\) with a uniform
remainder throughout \(T_n\asymp\log(1/\kappa_n)\), and construct an
honest connected interval attaining that expression.  In particular,
for every \(\lambda\in(0,\infty)\), derive and match the limiting
expected length along sequences satisfying
\(T_n/\log(1/\kappa_n)\to\lambda\).  The result must retain
effect-label testing, agree with the exact \(s=1\) common-outcome
plateau experiment, have a \(\lambda\downarrow0\) converse that
recovers Theorem 2's nonshrinking lower bound, with no matched
below-boundary value claimed without a separate attaining upper bound,
and agree in order with Theorems 1 and 3 as
\(\lambda\uparrow\infty\).  It must strictly extend the known-cutoff
honesty result of Noack and Rothe (2024) by evaluating connected
expected length uniformly over the unknown cutoff \(c_P\).
\end{openendedquestion}
\par\smallskip\noindent \textit{Justification.} The critical window is the only unresolved region where cutoff entropy and weak-denominator uncertainty remain of the same order, so evaluating it completes the frontier rather than juxtaposing regimes. \textit{Gap.} ArmstrongKolesar2018 evaluates fixed-level length through ordered moduli on convex fixed-index classes; the present experiment is a nonconvex union of translated causal ratio models with an effect-label coordinate. \textit{Consumer.} The resulting formula gives the signal threshold and interval length needed to plan, audit, and interpret fuzzy-RD studies whose operative cutoff is not recorded.

\begin{lemma}[lem:rectangular-jump-scan-localization]\label{lem:rectangular-jump-scan-localization}
Suppose \(s=1\), \(P\in\mathcal P_n\), and use the notation in the
definition of \(I_n^{\mathrm{SI}}\).  Write
\(\Delta_H^D(t)=E_P[\widehat\Delta_{j,H}^D(t)]\), which is independent
of \(j\), and define
\[
 \mathcal E_1=\left\{
 \max_{t\in\mathcal G_H}
 \left|\widehat\Delta_{1,H}^{D}(t)-\Delta_H^D(t)\right|
 \le\frac{\kappa_n}{16}\right\}.
\]
Then
\[
 P(\mathcal E_1^c)
 \le2N_H\exp\{-2^{-11}mH\kappa_n^2\},
 \qquad N_H\le1+32/H,
\]
every retained \(H\)-neighborhood is contained in
\([c_P-2H,c_P+2H]\) on \(\mathcal E_1\), and, for every
\(u\ge16/n\),
\[
 P\bigl(\{|\widehat c-c_P|>u\}\cap\mathcal E_1\bigr)
 \le9\exp\{-2^{-16}m\kappa_n^2u\}.
\]
\end{lemma}
\begin{proof}
Put \(f(x)=E_P[D\mid X=x]\).  By
\(\mathrm{ass:observable\mbox{-}holder\mbox{-}extensions}\) at
\(s=1\), the two one-sided extensions are \(M\)-Lipschitz.  Let
\(a=f(c_P-)\), let \(d=f(c_P+)-f(c_P-)=d_P\), and define
\[
 e(x)=\begin{cases}
 f(x)-a,&x<c_P,\\
 f(x)-a-d,&x\ge c_P.
 \end{cases}
\]
Both pieces vanish at \(c_P\).  Hence each is \(M\)-Lipschitz on its
side, and, if \(x<c_P<y\),
\[
 |e(y)-e(x)|\le M(y-c_P)+M(c_P-x)=M(y-x).
\]
Thus \(e\) is globally \(M\)-Lipschitz and
\[
 f(x)=a+d\mathbf 1\{x\ge c_P\}+e(x). \tag{1}
\]
The bounds \(\kappa_n\le1/8\) and \(M\ge2\) give
\(H\le2^{-16}<1/4\).  Together with
\(c_P,t\in[1/4,3/4]\), from
\(\mathrm{ass:interior\mbox{-}cutoff}\) and the scan definition, this
ensures that every \(H\)-window below is contained in the support
\([0,1]\); no boundary extrapolation is being used.
Uniformity of \(X\), from
\(\mathrm{ass:uniform\mbox{-}running\mbox{-}variable}\), and
\(\int_0^1K_0(u)\,du=1\) give, for all scan locations,
\[
 \Delta_H^D(t)=d\left(1-\frac{|t-c_P|}{H}\right)_++Q_e(t),
 \qquad
 Q_e(t)=\frac1H\int_0^H\{e(t+x)-e(t-x)\}\,dx. \tag{2}
\]
The Lipschitz property gives
\[
 |Q_e(t)|\le MH,
 \qquad |Q_e(t)-Q_e(q)|\le2M|t-q|. \tag{3}
\]
There is \(q\in\mathcal G_H\) with \(|q-c_P|\le H/64\), including
when \(c_P\) is a trimmed endpoint.  Since \(d\ge\kappa_n\) by
\(\mathrm{ass:local\mbox{-}compliance\mbox{-}strength}\) and
\(MH=\kappa_n/2^{12}\), (2)--(3) imply
\[
 \Delta_H^D(q)\ge\frac{63}{64}\kappa_n-\frac{\kappa_n}{2^{12}},
 \qquad
 \sup_{|t-c_P|\ge H}|\Delta_H^D(t)|
 \le\frac{\kappa_n}{2^{12}}. \tag{4}
\]

For fixed \(t\), one summand \(W_i(t)\) in the fixed-denominator
contrast has \(|W_i(t)|\le H^{-1}\),
\(E_P[W_i(t)^2]\le2H^{-1}\), and its centered version is bounded by
\(2H^{-1}\).  For \(|z|<3\), expansion of the exponential series and
\(k!\ge2\,3^{k-2}\) for \(k\ge2\) give
\[
 e^z\le1+z+\frac{z^2}{2(1-|z|/3)}.
\]
Applying this inequality conditionally to each centered summand,
multiplying the moment-generating-function bounds using
\(\mathrm{ass:iid\mbox{-}sampling}\), applying Markov's inequality,
and optimizing the exponential parameter yield
\[
 P\left\{\left|m^{-1}\sum_{i\in F_1}
       (W_i(t)-E_PW_i(t))\right|>x\right\}
 \le2\exp\left\{-\frac{mx^2}{4/H+4x/(3H)}\right\}. \tag{5}
\]
At \(x=\kappa_n/16\le1/128\), (5) is at most
\(2\exp\{-2^{-11}mH\kappa_n^2\}\).  A union bound over
\(N_H=1+\lfloor32/H\rfloor\le1+32/H\) grid points proves the first
display in the lemma.  On \(\mathcal E_1\), the empirical value at
\(q\) exceeds every empirical value at distance at least \(H\) from
\(c_P\) by more than \(\kappa_n/8\), by (4) and two deviations of
size \(\kappa_n/16\).  Every retained center is consequently within
\(H\) of \(c_P\), and its \(H\)-neighborhood is contained in
\([c_P-2H,c_P+2H]\).

For \(|t-c_P|\le H\), (2)--(3) imply
\[
 \Delta_H^D(c_P)-\Delta_H^D(t)
 \ge\left(\frac{\kappa_n}{H}-2M\right)|t-c_P|
 \ge\frac{\kappa_n}{2H}|t-c_P|. \tag{6}
\]
For \(H\le|t-c_P|\le2H\), the same difference is at least
\(\kappa_n-2MH\ge\kappa_n/2\).  On \(\mathcal E_1\), the set
\(B_H\) contains \(c_P\), because it contains the \(H\)-neighborhood
of its empirical coarse maximizer, which was just shown to lie within
\(H\) of \(c_P\).  Hence the componentwise grid contains a point
\(c_n\) with \(|c_n-c_P|\le1/n\).  The Lipschitz constant of the
right side of (2) is at most
\(d/H+2M\le3\kappa_n/H\).  Since \(u\ge16/n\), replacing \(c_P\)
by \(c_n\) consumes at most \(3\kappa_nu/(16H)\) of the separation
in (6).

We give the stopped exponential calculation for the local maximizer.
For \(0<r\le H\), moving \(t\) to the right by at most \(r\) changes
the rectangular weight only on the three boundary strips at
\(c_n-H,c_n,c_n+H\).  Their squared coefficients sum to six, so the
integrated squared weight is at most \(6r/H^2\), while the absolute
weight is at most \(2/H\).  Reveal the observations in the three
strips in increasing running-variable order, interleaving equal
relative positions.  Centering each newly revealed contribution gives
martingale increments bounded by \(4/H\); for \(2m\) observations the
predictable quadratic variation up to displacement \(R\) is at most
\(12mR/H^2\).  The exponential-series inequality used in (5) shows
that, for \(0<\eta H<3/4\),
\[
 \exp\left\{\eta S_k-
 \frac{\eta^2V_k}{2(1-4\eta/(3H))}\right\}
\]
is a nonnegative supermartingale.  Stopping it at the first crossing
and using Markov's inequality is legitimate because the scan grid is
finite.  On the annulus \(r_j=2^ju\le t-c_n<2r_j\), the deterministic
separation, including the grid error, is at least
\(\kappa_nr_j/(8H)\).  Thus the unnormalized sum of the \(2m\)
fold-two and fold-three contributions must cross
\(m\kappa_nr_j/(4H)\).  With \(\eta=\kappa_nH/2^7\),
\(\kappa_n\le1/8\), and \(V_k\le24mr_j/H^2\), the exponent is bounded
above by
\[
 -\frac{m\kappa_n^2r_j}{2^9}
 +\frac{12m\kappa_n^2r_j}
 {2^{14}(1-1/768)}
 \le-2^{-13}m\kappa_n^2r_j. \tag{7}
\]
The decreasing-order filtration gives the same bound on the left.
Once an annulus reaches distance \(H\), use the separation
\(\kappa_n/2\) and the full-window quadratic-variation bound; this is
no weaker than (7) after reducing its numerical exponent.  Summing the
dyadic annuli proves
\[
 P\bigl(\{|\widehat c-c_P|>u\}\cap\mathcal E_1\bigr)
 \le9\exp\{-2^{-16}m\kappa_n^2u\}.
\]
Indeed, if \(2^{-16}m\kappa_n^2u\le\log9\), the right side is at least
one; otherwise the two geometric series generated by (7) are bounded
by the displayed right side.  For \(u>2H\), the event is empty by the
already proved containment.  This proves every assertion of the lemma.
\end{proof}
\par\smallskip\noindent \textit{Justification.} This lemma supplies the population triangular profile, the global scan bound, and the stopped local-refinement inequality for the actual rectangular sample map. \textit{Gap.} The exponential \(m\kappa_n^2u\) tail is derived with displayed constants rather than borrowed from fixed-jump asymptotics. \textit{Consumer.} It is the localization input for the Lipschitz honesty and expected-length proof.

\begin{theorem}[thm:lipschitz-strong-signal-oracle-attainment]\label{thm:lipschitz-strong-signal-oracle-attainment}
On the Lipschitz face \(s=1\), there are finite constants
\(C_1,C_2\), depending only on \(M,v,\alpha\), such that, for all
sufficiently large \(n\), if
\[
 T_n=n\kappa_n^3\ge C_1\log(e/\kappa_n),
\]
then \(I_n^{\mathrm{SI},1}\) is uniformly honest over
\(\mathcal P_n(1,M,v,\kappa_n)\) and
\[
 \sup_{P\in\mathcal P_n}E_P[|I_n^{\mathrm{SI},1}|]
 \le C_2\min\{1,\rho_n/\kappa_n\}.
\]
Consequently
\(L_n(\kappa_n)\asymp L_n^{\mathrm{known}}(\kappa_n)\) in this
region, uniformly over sequences \(\kappa_n\in(0,1/8]\), where at
\(s=1\) one has \(\rho_n=n^{-1/3}\).
\end{theorem}
\begin{proof}
On the asserted face,
\[
 H=\frac{\kappa_n}{2^{12}M},\qquad
 b=m^{-1/3},\qquad
 T_n=n\kappa_n^3,\qquad
 \rho_n=n^{-1/3}. \tag{8}
\]
Put
\[
 L_\alpha=\log(32/\alpha),\qquad A_\alpha=2^{16}L_\alpha,
 \qquad
 B_{M,\alpha}=1+\log\!\left(
   \frac{16(1+2^{17}M)}{\alpha}\right),
\]
and define the finite constants
\[
 C_r=8M+16A_\alpha+1+16M(A_\alpha+1)+8\sqrt{L_\alpha},
\]
\[
 c_3=2^{-24}M^{-1},\qquad c_2=2^{-22},
\]
\[
 C_F=2(3ec_3)^{-1/3}+12(2ec_2)^{-1/2},
 \qquad C_2=16C_r+4C_F,
\]
\[
 C_1=\max\left\{
 2^{27}MB_{M,\alpha},\;
 2^{15}M(8A_\alpha+1),\;
 8(8C_r)^3,\;
 2^{30}ML_\alpha
 \right\}. \tag{9}
\]
These constants are deliberately conservative.  The \(m=0\) branch
in the definition makes the rule total for every \(n\); for all
sufficiently large \(n\), which is the range of the theorem,
\(m\ge n/8\ge1\), \(b\le1/8\), and the nontrivial branch is available.
Since \(\widehat c\in[1/4,3/4]\), the inequality \(b\le1/8\) also
places every fold-four window inside \([0,1]\).

Since \(N_H\le1+2^{17}M/\kappa_n\),
\[
 \log(16N_H/\alpha)
 \le B_{M,\alpha}\log(e/\kappa_n). \tag{10}
\]
Moreover,
\[
 2^{-11}mH\kappa_n^2
 \ge\frac{T_n}{2^{26}M}. \tag{11}
\]
The first term in (9), the signal condition, and
\(\mathrm{lem:rectangular\mbox{-}jump\mbox{-}scan\mbox{-}localization}\)
therefore give
\[
 P(\mathcal E_1^c)\le\alpha/8. \tag{12}
\]
The same lemma, at
\(R_c=A_\alpha/(m\kappa_n^2)+1/n\), yields
\[
 P\bigl(\{|\widehat c-c_P|>R_c\}\cap\mathcal E_1\bigr)
 \le9e^{-L_\alpha}=\frac{9\alpha}{32}, \tag{13}
\]
because \(R_c\ge16/n\), using \(m\le n/4\),
\(\kappa_n\le1/8\), and \(A_\alpha>1\).

We next verify the deterministic fallbacks.  Let
\(q=\kappa_n/b\).  Since \(b^{-3}=m\) and \(m\ge n/8\),
\[
 q^3=m\kappa_n^3\ge T_n/8. \tag{14}
\]
The signal condition and (9) imply \(q\ge1\).  Direct substitution,
using \(n\ge4m\), gives
\[
 r_n\le C_rb. \tag{15}
\]
Indeed, after division by \(b\), the four terms in \(r_n\) are bounded
respectively by \(8M\), \(16A_\alpha+1\),
\(16M(A_\alpha+1)\), and \(8\sqrt{L_\alpha}\).  Equations (9),
(14), and (15) imply \(r_n\le\kappa_n/8\).  Also
\[
 \frac{R_c}{H}
 =2^{12}M\left\{
 \frac{A_\alpha}{m\kappa_n^3}+\frac1{n\kappa_n}
 \right\}
 \le\frac{2^{12}M(8A_\alpha+1)}{T_n}
 \le\frac18, \tag{16}
\]
where \(n\kappa_n\ge n\kappa_n^3=T_n\).  Thus every deterministic
fallback check passes under the stated signal condition for all
sufficiently large \(n\).

Fix \(P\in\mathcal P_n\) and write \(\theta=\theta_P\).  By
\(\mathrm{prop:wald\mbox{-}late\mbox{-}identification}\),
\(y_P-\theta d_P=0\), and \(|\theta|\le1\) because
\(\theta\in\Theta=[-1,1]\).  The conditional regression of
\(U=Y-\theta D\), separately on either side of \(c_P\), is therefore
\(2M\)-Lipschitz and has zero level jump.  Gluing the two pieces at
their common one-sided value makes that regression globally
\(2M\)-Lipschitz.  Conditional on folds one through three,
\(\widehat c\) is fixed and fold four is independent by
\(\mathrm{ass:iid\mbox{-}sampling}\).  On
\(|\widehat c-c_P|\le R_c\), the rectangular moment identity gives
\[
 \left|E_P\left[
 \widehat y-\theta\widehat d\mid F_1,F_2,F_3\right]\right|
 \le2Mb. \tag{17}
\]
For treatment, applying (2)--(3) from the localization proof at
bandwidth \(b\), and using \(d_P\le2\kappa_n\), gives
\[
 \left|E_P[\widehat d\mid F_1,F_2,F_3]-d_P\right|
 \le Mb+\frac{2\kappa_nR_c}{b}+2MR_c. \tag{18}
\]
Every division is legitimate because \(m\ge1\), \(b>0\), and
\(\kappa_n>0\).  Moreover, the nonnegative term
\(16\kappa_nR_c/b\) in \(r_n\), together with
\(r_n\le\kappa_n/8\), implies \(R_c/b\le1/128\), so the shifted
cutoff lies inside the kernel window.

For completeness, the fold-four concentration follows directly from
the exponential-series inequality used in the localization proof.
The summand for \(U\) is bounded by \(2/b\), is supported on a set of
\(X\)-probability \(2b\), and has second moment at most \(8/b\); the
corresponding bounds for \(D\) are \(1/b\) and \(2/b\).  Multiplication
of the conditional moment-generating functions and Markov's inequality
therefore show, for all sufficiently large \(n\), that
\[
 \left|\widehat y-\theta\widehat d
 -E_P[\widehat y-\theta\widehat d\mid F_1,F_2,F_3]\right|
 \le8\sqrt{\frac{L_\alpha}{mb}} \tag{19}
\]
except on an event of probability at most \(\alpha/8\).  The linear
term in the Bernstein denominator is absorbed because
\(mb=m^{2/3}\to\infty\).  Combining (17), (19), and the definition of
\(r_n\) yields
\[
 |\widehat y-\theta_P\widehat d|\le2r_n. \tag{20}
\]
If \(\widehat d<\kappa_n/2\), the rule returns \(\Theta\), which
contains \(\theta_P\).  Otherwise (20) puts \(\theta_P\) in the
inverted set.  By (12), (13), and (19), the total failure probability
is at most
\[
 \frac\alpha8+\frac{9\alpha}{32}+\frac\alpha8
 =\frac{17\alpha}{32}<\alpha.
\]
The bounds are uniform in \(P\), proving honesty.

When \(\widehat d\ge\kappa_n/2\), the accepted labels form an interval
of length at most
\[
 \frac{4r_n}{\widehat d}\le\frac{8r_n}{\kappa_n}. \tag{21}
\]
On \(|\widehat c-c_P|\le b/8\), (18),
\(Mb\le r_n/8\le\kappa_n/64\),
\(2MR_c\le r_n/8\le\kappa_n/64\), and a fold-four treatment
fluctuation at most \(\kappa_n/8\) imply
\(\widehat d>\kappa_n/2\).  Hence the full-length fallback event is
contained in a coarse-scan failure, a localization error exceeding
\(b/8\), or a fold-four treatment deviation exceeding
\(\kappa_n/8\).  From (10)--(14), the localization lemma, and the same
fold-four exponential calculation,
\[
 P\{\widehat d<\kappa_n/2\}
 \le2e^{-c_3q^3}+12e^{-c_2q^2}. \tag{22}
\]
Here the first term follows because the first term of (9) gives
\(\log N_H\le q^3/(2^{24}M)\); the localization exponent at \(b/8\)
is at least \(2^{-19}q^2\); and the treatment Bernstein exponent is
larger than \(2^{-22}q^2\).  Thus every constant in (22) follows from
the preceding displayed inequalities.

The elementary one-variable optimizations
\[
 \sup_{q\ge1}qe^{-c_3q^3}\le(3ec_3)^{-1/3},
 \qquad
 \sup_{q\ge1}qe^{-c_2q^2}\le(2ec_2)^{-1/2}
\]
turn (22) into
\[
 P\{\widehat d<\kappa_n/2\}
 \le C_F\frac{b}{\kappa_n}. \tag{23}
\]
Because \(m\ge n/8\), \(b\le2\rho_n\).  Taking expectations in
(21), using the maximal length two on the fallback, and applying
(15) and (23) give
\[
 \sup_{P\in\mathcal P_n}E_P|I_n^{\mathrm{SI},1}|
 \le\frac{8C_rb}{\kappa_n}+\frac{2C_Fb}{\kappa_n}
 \le C_2\frac{\rho_n}{\kappa_n}. \tag{24}
\]
Equations (9) and (14) imply \(\rho_n/\kappa_n<1\) in the asserted
region.  Thus (24) has the claimed minimum form.

Finally, every unknown-cutoff interval rule induces a known-cutoff rule
that ignores its supplied value of \(c_P\), so
\(L_n^{\mathrm{known}}(\kappa_n)\le L_n(\kappa_n)\).  The already
proved \(\mathrm{thm:known\mbox{-}cutoff\mbox{-}benchmark}\), the
definition of \(L_n\), and (24) imply
\[
 c_0\min\{1,\rho_n/\kappa_n\}
 \le L_n^{\mathrm{known}}(\kappa_n)
 \le L_n(\kappa_n)
 \le C_2\min\{1,\rho_n/\kappa_n\}.
\]
This proves the Lipschitz oracle-order assertion.
\end{proof}
\par\smallskip\noindent \textit{Justification.} This theorem is the proved strong-signal result: an explicit total interval attains the oracle connected-length order on the Lipschitz face. \textit{Gap.} It extends fixed-jump cutoff substitution to a vanishing compliance jump with uniform honest expected length, but only under the existing \(s=1\) observable class. \textit{Consumer.} It gives implementable bandwidths, fallbacks, and a signal threshold for Lipschitz fuzzy-RD studies with an unrecorded cutoff.

\begin{lemma}[lem:profile-cutoff-localization]\label{lem:profile-cutoff-localization}
Fix \(1<s\le2\), \(M\ge2\), and \(v\in(0,1/64)\).  The asserted
uniform localization inequality for
\((\widehat c_{\mathrm{prof}},\mathsf F_{\mathrm{prof}})\) is false on
\(\mathcal P_n(s,M,v,\kappa_n)\).  More precisely, for every
\[
 0<\varepsilon<\frac{1}{2(2+1/s)},
 \qquad \kappa_n=\frac18n^{-\varepsilon},
 \qquad H_n=\left(\frac{\kappa_n}{2^{12}M}\right)^{1/s},
 \qquad m=\lfloor n/4\rfloor,
\]
there is, for all sufficiently large \(n\), a law
\(P_n\in\mathcal P_n(s,M,v,\kappa_n)\) with \(c_{P_n}=1/2\) such that
\[
 \frac{T_n}{\log(e/\kappa_n)}\longrightarrow\infty,
 \qquad \frac{16}{n}\le\frac{H_n}{8},
\]
and
\[
 P_n\{\mathsf F_{\mathrm{prof}}=1\}
 +P_n\!\left(\left\{|\widehat c_{\mathrm{prof}}-c_{P_n}|>
                    H_n/8\right\}
       \cap\{\mathsf F_{\mathrm{prof}}=0\}\right)
 \longrightarrow1.
\]
Consequently, there are no positive constants
\(A_{\mathrm{prof}},a_{\mathrm{prof}},C_{\mathrm{prof}}\), depending
only on \(s,M\), such that, simultaneously for every
\(P\in\mathcal P_n(s,M,v,\kappa_n)\), every index satisfying
\(T_n\ge C_{\mathrm{prof}}\log(e/\kappa_n)\), and every
\(16/n\le u\le H_n\),
\[
 P\{\mathsf F_{\mathrm{prof}}=1\}
 +P\!\left(\{|\widehat c_{\mathrm{prof}}-c_P|>u\}
       \cap\{\mathsf F_{\mathrm{prof}}=0\}\right)
 \le A_{\mathrm{prof}}\exp\{-a_{\mathrm{prof}}m\kappa_n^2u\}.
\]
\end{lemma}
\begin{proof}
Fix \(s\in(1,2]\), \(M\ge2\), \(v\in(0,1/64)\), and choose
\[
 0<\varepsilon<\frac{1}{2(2+1/s)},
 \qquad
 \kappa_n=\frac18 n^{-\varepsilon},
 \qquad
 H_n=\left(\frac{\kappa_n}{2^{12}M}\right)^{1/s},
 \qquad c=\frac12,
 \qquad m=\lfloor n/4\rfloor.
 \tag{1}
\]
Write \(q_n=\kappa_n/H_n\).  Since \(s>1\),
\[
 q_n=(2^{12}M)^{1/s}\kappa_n^{\,1-1/s}\longrightarrow0.
 \tag{2}
\]

We first construct the full-data law.  Let \(X\) be uniform on
\([0,1]\).  Conditional on \(X=x\), assign the monotone response type
\(R\) the probabilities
\[
 P_n(R=A\mid X=x)=\frac14,
 \qquad
 P_n(R=C\mid X=x)=w_n(x),
 \qquad
 P_n(R=N\mid X=x)=\frac34-w_n(x),
 \tag{3}
\]
where
\[
 w_n(x)=
 \begin{cases}
  \kappa_n,&x<c,\\
  \kappa_n+q_n(x-c),&x\ge c.
 \end{cases}
 \tag{4}
\]
By (2), \(0\le w_n(x)\le\kappa_n+q_n/2<1/4\) for all sufficiently
large \(n\); hence the three masses in (3) are nonnegative and sum to
one.  Given \((X,R)\), let \(U\) be Bernoulli with success probability
\(1/2\), independently of \((X,R)\), and set
\[
 Y(0)=Y(1)=U.
 \tag{5}
\]
Use the frozen threshold assignment \(Z=\mathbf 1\{X\ge c\}\), the
frozen monotone treatment rule \(D=D(Z)\), and independent copies
across units.

We verify every condition defining the model class.  From (3)--(4),
\[
 E_{P_n}[D(0)\mid X=x]=\frac14,
 \qquad
 E_{P_n}[D(1)\mid X=x]=\frac14+w_n(x).
 \tag{6}
\]
Both functions in (6) are continuous at \(c\), because the two
one-sided values of \(w_n\) there are \(\kappa_n\).  By (5),
\(E_{P_n}[Y(D(z))\mid X=x]=1/2\) for both \(z\in\{0,1\}\).  Thus
\(\mathrm{ass:potential\mbox{-}reduced\mbox{-}form\mbox{-}continuity}\)
holds.  The observable treatment regressions are
\[
 \mu_{D,P_n}^{-}(x)=\frac14,
 \qquad
 \mu_{D,P_n}^{+}(x)=\frac14+\kappa_n+q_n(x-c).
 \tag{7}
\]
The left side has a constant extension, and the displayed affine
formula extends the right side to \([0,1]\).  For all sufficiently
large \(n\), these extensions take values in \([0,1]\), have uniformly
bounded suprema, have derivative suprema at most \(q_n\), and have
zero derivative H\"older seminorm.  Equation (2) and \(M\ge2\)
therefore place them in the radius-\(M\), order-\(s\) H\"older ball
required by
\(\mathrm{ass:observable\mbox{-}holder\mbox{-}extensions}\).  The two
observable outcome regressions are the constant \(1/2\), so they obey
the same condition.  Moreover,
\[
 d_{P_n}=\mu_{D,P_n}^{+}(c)-\mu_{D,P_n}^{-}(c)=\kappa_n,
 \tag{8}
\]
which verifies \(\mathrm{ass:local\mbox{-}compliance\mbox{-}strength}\).
Finally, (5) and the independence of \(U\) give, on both sides
\(\ell\),
\[
 \operatorname{Var}_{P_n}(Y\mid X=x,\ell)=\frac14>v,
 \qquad
 \operatorname{Cov}_{P_n}(Y,D\mid X=x,\ell)=0.
 \tag{9}
\]
Thus the variance floor and correlation-separation conditions hold.
The sampling is independent, the design is uniform, and \(c=1/2\) is
interior.  Equations (3), (5)--(9), together with those three facts,
prove
\[
 P_n\in\mathcal P_n(s,M,v,\kappa_n)
 \quad\text{for all sufficiently large }n.
 \tag{10}
\]

We next evaluate exactly the population criterion induced by the
moving-window profile.  For a candidate \(t\) satisfying
\(|t-c|\le H_n\), put
\[
 z=\frac{x-t}{H_n},
 \qquad
 r=\frac{c-t}{H_n},
 \qquad
 f_r(z)=\mathbf 1\{z\ge r\}(1+z-r).
 \tag{11}
\]
After subtracting the common affine function \(1/4\), (7) is
\(\kappa_nf_r(z)\) on \([t-H_n,t+H_n]\).  Write
\(\mu_{D,P_n}(x)=E_{P_n}[D\mid X=x]\) for the observable conditional
treatment mean, equal to the appropriate side of (7).  For every measurable
prediction \(g\), conditional variance decomposition gives
\[
E_{P_n}[\{D-g(X)\}^2\mid X=x]
=\operatorname{Var}_{P_n}(D\mid X=x)
 +\{\mu_{D,P_n}(x)-g(x)\}^2.
\tag{12}
\]
The variance term in (12) is identical for the common-affine and
piecewise-affine models, so it cancels from their risk difference.
Let \(\mathcal A_{n,t}^0\) and \(\mathcal A_{n,t}^1\) denote,
respectively, the constrained common-affine and constrained
piecewise-affine classes used by the frozen construction at \(t\).  We
define, before using either object,
\[
 R_{n,t}(g)=E_{P_n}\!\left[
   \mathbf 1\{|X-t|\le H_n\}\{D-g(X)\}^2\right],
 \qquad
 Q_n(t)=\inf_{g\in\mathcal A_{n,t}^0}R_{n,t}(g)
       -\inf_{g\in\mathcal A_{n,t}^1}R_{n,t}(g).
\]

Here is the complete projection calculation.  With the basis
\(e(z)=(1,z)^\top\), the common, left-flank, and right-flank Gram
matrices and their inverses are
\[
 G_0=\int_{-1}^{1}e(z)e(z)^\top dz
   =\begin{pmatrix}2&0\\0&2/3\end{pmatrix},
 \qquad
 G_0^{-1}=\begin{pmatrix}1/2&0\\0&3/2\end{pmatrix},
 \tag{13}
\]
\[
 G_- =\int_{-1}^{0}e(z)e(z)^\top dz
   =\begin{pmatrix}1&-1/2\\-1/2&1/3\end{pmatrix},
 \qquad
 G_-^{-1}=\begin{pmatrix}4&6\\6&12\end{pmatrix},
 \tag{14}
\]
and
\[
 G_+ =\int_{0}^{1}e(z)e(z)^\top dz
   =\begin{pmatrix}1&1/2\\1/2&1/3\end{pmatrix},
 \qquad
 G_+^{-1}=\begin{pmatrix}4&-6\\-6&12\end{pmatrix}.
 \tag{15}
\]
The determinants are \(4/3\) for \(G_0\) and \(1/12\) for each of
\(G_-\) and \(G_+\).  They are therefore positive definite and full
rank, so every inversion in (13)--(15) is permitted.
If \(0\le r\le1\), the left moment is zero and both the common and
right moments equal
\[
 u(r)=\int_r^1 e(z)(1+z-r)\,dz
 =\begin{pmatrix}
 (r-3)(r-1)/2\\[2pt]
 (r-1)(r^2-2r-5)/6
 \end{pmatrix}.
 \tag{16}
\]
If \(-1\le r\le0\), the left, right, and common moments are,
respectively,
\[
 u_-(r)=\int_r^0 e(z)(1+z-r)\,dz
 =\begin{pmatrix}r(r-2)/2\\ r^2(r-3)/6\end{pmatrix},
 \tag{17}
\]
\[
 u_+(r)=\int_0^1 e(z)(1+z-r)\,dz
 =\begin{pmatrix}-(2r-3)/2\\ -(3r-5)/6\end{pmatrix},
 \qquad
 u_0(r)=u_-(r)+u_+(r)
 =\begin{pmatrix}
 (r-3)(r-1)/2\\[2pt]
 (r-1)(r^2-2r-5)/6
 \end{pmatrix}.
 \tag{18}
\]
For a regression function with moment vector \(u\) and Gram matrix
\(G\), the squared norm of its least-squares projection is
\(u^\top G^{-1}u\).  Let \(\overline Q_n(t)\) denote common risk minus
piecewise risk when the affine coefficient constraints are omitted.
After the change of variables \(x=t+H_nz\),
\[
 \overline Q_n(t)=H_n\kappa_n^2q(r),
 \tag{19}
\]
where (13)--(18) give, without an omitted remainder,
\[
 q(r)=
 \begin{cases}
 u_-(r)^\top G_-^{-1}u_-(r)
 +u_+(r)^\top G_+^{-1}u_+(r)
 -u_0(r)^\top G_0^{-1}u_0(r),&-1\le r\le0,\\[3pt]
 u(r)^\top G_+^{-1}u(r)-u(r)^\top G_0^{-1}u(r),&0\le r\le1,
 \end{cases}
 \tag{20}
\]
\[
 q(r)=
 \begin{cases}
 \displaystyle
 \frac{(1+r)^2(7r^4-32r^3+27r^2+22r+4)}{24},&-1\le r\le0,\\[7pt]
 \displaystyle
 \frac{(1-r)^2(7r^4-52r^3+99r^2-10r+4)}{24},&0\le r\le1.
 \end{cases}
 \tag{21}
\]
This also verifies directly that the population projection is the
object being evaluated, rather than assuming the desired profile
separation.  On the compact set of \(r\)'s used below, the coefficients
obtained from (13)--(18) are bounded numerical multiples of
\(\kappa_n\) in the \(z\) parametrization.  In the original
\(x\) parametrization, their slopes are therefore \(O(q_n)\).  Their
local-coordinate intercepts are \(1/4+O(\kappa_n)\); if an affine
function is instead written with a global \(x\)-intercept, that
intercept is \(1/4+O(q_n)\).  By (2), the coefficients in either
parametrization are strictly inside the compact constraints in
\(\mathrm{def:profile\mbox{-}cutoff\mbox{-}center}\) for all large
\(n\).  Hence \(Q_n(t)=\overline Q_n(t)\), and (19)--(21) hold with
\(Q_n\) in place of \(\overline Q_n\), at every point used in the
comparison.

The exact formula has a strict displacement gain:
\[
 q(0)=\frac16,
 \qquad
 q\left(\frac25\right)=\frac{5949}{31250},
 \qquad
 \Delta:=q\left(\frac25\right)-q(0)
       =\frac{1111}{46875}>0.
 \tag{22}
\]
Moreover, \(q(r)\le1/6\) for \(|r|\le1/8\).  On the negative half,
differentiating the first polynomial in (21) gives
\[
 q'(r)=\frac{(1+r)(7r^4-22r^3+2r^2+20r+5)}4>0,
 \qquad -\frac18\le r\le0,
 \tag{23}
\]
because the second factor is at least
\(5+20r\ge5/2\).  On the positive half, exact factorization gives
\[
 \frac16-q(r)
 =-\frac{r(r^2-4r+1)(7r^3-38r^2+51r-18)}{24}\ge0,
 \qquad 0\le r\le\frac18.
 \tag{24}
\]
Indeed, \(r^2-4r+1>0\) on this interval, whereas
\(7r^3-38r^2+51r-18\le7/512+51/8-18<0\).  This proves the asserted
local bound.

Let \(t_n^\circ=c-(2/5)H_n\).  This point is in \([1/4,3/4]\) because
\(H_n<1/4\).  The deterministic grid in the frozen construction has
mesh at most
\(\eta_n=\min\{n^{-1},m^{-1/3}H_n/64\}\), so it contains a point
\(\bar t_n\) such that
\[
 |\bar t_n-t_n^\circ|\le\eta_n,
 \qquad
 \frac{\eta_n}{H_n}\le\frac{m^{-1/3}}{64}\longrightarrow0.
 \tag{25}
\]
Continuity of the polynomial in (21), together with (22)--(25), yields
eventually
\[
 Q_n(\bar t_n)
 \ge H_n\kappa_n^2\left(\frac16+\frac{3\Delta}{4}\right),
 \qquad
 \sup_{|t-c|\le H_n/8}Q_n(t)
 \le\frac{H_n\kappa_n^2}{6}.
 \tag{26}
\]

It remains to show that random fitting cannot reverse the fixed gap in
(26).  We give the empirical comparison rather than assume a maximal
inequality.  Every localized squared-loss function used on the first
three folds is indexed by one interval center and at most four affine
coefficients in the fixed compact coefficient boxes of the
construction.  Since \(D\in\{0,1\}\), these losses have an envelope
\(B_M<\infty\) depending only on \(M\).  Put
\(\delta_m=m^{-3}\), take a deterministic \(\delta_m\)-mesh of the
cutoff-coordinate interval, and put a deterministic
\(\delta_m\)-net on every coefficient box.  There are at most
\(C_Mm^{C_M}\) representatives.  If \(t\) is replaced by its nearest
mesh point, every changed support or flank indicator lies in one of a
fixed collection of \(O(m^3)\) boundary bands, each having Lebesgue
length at most \(C\delta_m\).  Uniformity of \(X\) gives probability at
most \(C\delta_m\) to each band.  For a fold of size \(m\), write
\(N_B\) for its number of observations in a boundary band \(B\).  The
union bound and
\[
 P_n\{N_B\ge4\}\le {m\choose4}(C\delta_m)^4
\]
show that the probability some such band contains four observations
is \(O(m^3m^4\delta_m^4)=O(m^{-5})\).  Off this event, replacing \(t\)
by its mesh representative changes an empirical risk by
\(O_M(m^{-1})\).  If the affine slope is represented in the normalized
flank coordinate, its loss on an unchanged support varies by at most
\(O_M(\delta_m/H_n)\); in an unnormalized coordinate the smaller bound
\(O_M(\delta_m)\) applies.  Coefficient discretization contributes
\(O_M(\delta_m)\).  The corresponding population-risk errors satisfy
the same bounds, in addition to the \(O_M(\delta_m)\) boundary-mass
term, by the uniform design.  From (1),
\(m^{-3}/H_n=o(m^{-1})\).  Thus the deterministic finite net
approximates empirical and population risks uniformly by
\(O_M(m^{-1})\) on an event whose probability tends to one.

For completeness, the required finite-net concentration follows
directly from boundedness.  If \(W\) is centered and
\(|W|\le B_M\), convexity of the exponential on
\([-B_M,B_M]\) gives
\[
 E e^{aW}\le\cosh(aB_M)\le e^{a^2B_M^2/2}.
 \tag{27}
\]
The first inequality is the expectation of the chord bound for
\(e^{aw}\) on \([-B_M,B_M]\), using \(E[W]=0\).  The second follows
term by term from
\((2k)!\ge2^k k!\) in the power series for \(\cosh\).
Exponential Markov inequality applied to independent copies and then
optimized at \(a=t/B_M^2\) therefore gives
\[
 P\left(\left|m^{-1}\sum_{i=1}^mW_i\right|>t\right)
 \le2\exp\left\{-\frac{mt^2}{2B_M^2}\right\}.
 \tag{28}
\]
A union bound over the polynomial net, followed by the preceding
approximation, proves on each of the first three folds
\[
 \sup_{t,g}\left|
  \frac1m\sum_{i\in F_j}\mathbf 1\{|X_i-t|\le H_n\}
       \{D_i-g(X_i)\}^2
  -E_{P_n}\!\left[\mathbf 1\{|X-t|\le H_n\}
       \{D-g(X)\}^2\right]
 \right|
 =O_{P_n}\!\left(\sqrt{\frac{\log n}{n}}\right).
 \tag{29}
\]
The same construction with separate left- and right-flank indicators
covers the piecewise class.  Because the supremum in (29) ranges over
every admissible coefficient vector, it also covers a fitted function
trained on an independent fold.

Let \(e_n\) denote a common random upper bound in (29), defined over
the full deterministic loss classes.  Let \(\mathcal G_n\) be the
event that all empirical Gram checks in the frozen construction pass.
Only on \(\mathcal G_n\) are the fitted scores and their maximizer
formed.  On that event, uniform empirical-risk minimization implies,
for each model \(k\in\{0,1\}\), training fold \(j\), and candidate
\(t\),
\[
 0\le R_{n,t}(\widehat g_{j,t}^{k})
       -\inf_{g\in\mathcal A_{n,t}^k}R_{n,t}(g)\le2e_n.
 \tag{30}
\]
Applying (29) on the held-out fold and using (30) for the common and
piecewise fits gives, on \(\mathcal G_n\),
\[
 \sup_{t\in\mathcal T_n}|\widehat Q(t)-Q_n(t)|
 =O_{P_n}\!\left(\sqrt{\frac{\log n}{n}}\right).
 \tag{31}
\]
On \(\mathcal G_n^c\), the construction stops
before defining a score or a maximizer and returns
\(\mathsf F_{\mathrm{prof}}=1\); no undefined maximizer is invoked.
The empirical-process event underlying (29) is defined directly over
the full deterministic loss classes and hence is meaningful whether
or not \(\mathcal G_n\) occurs.

By (1),
\[
 \frac{\sqrt{\log n/n}}{H_n\kappa_n^2}
 =O\!\left(n^{-1/2+\varepsilon(2+1/s)}\sqrt{\log n}\right)
 \longrightarrow0.
 \tag{32}
\]
Let \(\mathcal E_n\) be the event, defined through the uniform
empirical-risk deviations over the deterministic classes, on which
those deviations are small enough that, whenever \(\mathcal G_n\)
occurs, the supremum in (31) is smaller than
\(H_n\kappa_n^2\Delta/4\).  Equations (29) and (32) give
\(P_n(\mathcal E_n)\to1\), whether or not the Gram checks pass.  On
\(\mathcal E_n\cap\mathcal G_n\), (26) shows that the score at
\(\bar t_n\) is strictly larger than the score at every grid point in
\([c-H_n/8,c+H_n/8]\).  The smallest maximizer \(\widetilde c\), which
is defined on \(\mathcal G_n\), must therefore satisfy
\[
 |\widetilde c-c|>H_n/8
 \quad\text{on }\mathcal E_n\cap\mathcal G_n.
 \tag{33}
\]
If \(\mathcal G_n^c\) occurs, the first probability in the target
display applies.  If \(\mathcal E_n\cap\mathcal G_n\) occurs, then the
later fitted-jump check either fails, in which case it again returns
\(\mathsf F_{\mathrm{prof}}=1\), or it passes, in which case
\(\widehat c_{\mathrm{prof}}=\widetilde c\) and (33) places the center
outside the stated radius.  The two events counted in the next display
are disjoint, and their union contains \(\mathcal E_n\).  Consequently,
\[
 P_n\{\mathsf F_{\mathrm{prof}}=1\}
 +P_n\!\left(\{|\widehat c_{\mathrm{prof}}-c|>H_n/8\}
       \cap\{\mathsf F_{\mathrm{prof}}=0\}\right)
 \ge P_n(\mathcal E_n)\longrightarrow1.
 \tag{34}
\]

Finally,
\[
 T_n=n\kappa_n^{2+1/s}
 =8^{-(2+1/s)}n^{\,1-\varepsilon(2+1/s)},
 \tag{35}
\]
so \(T_n/\log(e/\kappa_n)\to\infty\).  Also
\(nH_n\to\infty\), and therefore \(16/n\le H_n/8\) eventually.  For
arbitrary fixed positive
\(A_{\mathrm{prof}},a_{\mathrm{prof}},C_{\mathrm{prof}}\), the signal
condition with \(C_{\mathrm{prof}}\) eventually holds, whereas the
claimed right side at \(u_n=H_n/8\) is
\[
 A_{\mathrm{prof}}
 \exp\left\{-\frac{a_{\mathrm{prof}}m\kappa_n^2H_n}{8}\right\}
 =A_{\mathrm{prof}}
 \exp\left\{-\frac{a_{\mathrm{prof}}m}{8n(2^{12}M)^{1/s}}T_n\right\}
 \longrightarrow0,
 \tag{36}
\]
because \(m/n\to1/4\) and (35) diverges.  Equations (34) and (36)
contradict the frozen positive inequality for every choice of its three
constants and prove the stated negative result.
\end{proof}
\par\smallskip\noindent \textit{Justification.} This counterexample is load-bearing because it prevents the higher-smoothness argument from using a false localization property of the frozen moving-window profile. \textit{Gap.} The obstruction is deterministic rather than a missing maximal inequality: within the unchanged class, the population profile can strictly prefer a displaced split to every split near the true cutoff. \textit{Consumer.} The higher-smoothness detectable-region OEQ uses this result to exclude the frozen profile as an attaining witness and to require a genuinely different estimator or a negative frontier conclusion.

\begin{openendedquestion}[oeq:smooth-detectable-region-bracket]\label{oeq:smooth-detectable-region-bracket}
For \(1<s\le2\), determine under the unchanged class
\(\mathcal P_n(s,M,v,\kappa_n)\) whether a genuinely different total
estimator or cutoff profile, together with same-class slope-kink lower
pairs, establishes the detectable-region frontier
\[
 L_n(\kappa_n)\asymp
 \min\left\{1,\frac{\rho_n}{\kappa_n}
                  +\frac1{n\kappa_n^3}\right\}
 \qquad\text{when}\qquad
 T_n\gtrsim\log(e/\kappa_n).
\]
The frozen common-versus-piecewise affine moving-window profile in the
\(1<s\le2\) branch of \(I_n^{\mathrm{SI}}\) cannot establish this upper
bracket through the proposed localization route:
the counterexample in
\(\mathrm{lem:profile\mbox{-}cutoff\mbox{-}localization}\) has
\(T_n/\log(e/\kappa_n)\to\infty\), yet the profile falls back or selects
a center farther than \(H_n/8\), where
\(H_n=(\kappa_n/(2^{12}M))^{1/s}\), with probability tending to one
because its population score is larger at a displaced split.  A positive
resolution must therefore construct and analyze a genuinely different
estimator or profile; alternatively, a negative resolution must prove
that the proposed frontier itself is unattainable.  Either resolution
must give explicit constants for uniform honesty and expected length,
derive a same-class lower construction for the term
\(1/(n\kappa_n^3)\), and decide whether oracle order is both sufficient
and necessary on the face \(n\kappa_n^2\rho_n\gtrsim1\).  No
cross-side derivative-matching assumption may be added, and the frozen
profile branch may not be reused as an attaining witness.
\end{openendedquestion}
\par\smallskip\noindent \textit{Justification.} This question now carries the genuine higher-smoothness frontier problem after the frozen moving-window profile has been refuted on the unchanged class. \textit{Gap.} A positive answer requires a new total estimator or cutoff profile with a valid uniform upper proof, while a matched frontier also requires a same-class slope-kink lower pair; the frozen profile cannot supply the former through its proposed localization route. \textit{Consumer.} Its answer will determine whether cutoff uncertainty contributes the candidate term beyond the known-cutoff rate and, if so, which replacement procedure attains it.

\begin{proposition}[prop:outcome-invariant-witness-high-separation]\label{prop:outcome-invariant-witness-high-separation}
Suppose \(h_n\le 1/32\), with \(a_\star=1-v\) and \(b=1-v/2\).  Then \(J_n\ge2\), the cutoff locations in \(\mathcal W_n\) lie in \([1/4,3/4]\), and the supports of distinct compliance profiles are disjoint up to Lebesgue-null endpoints.  The common potential-outcome table satisfies
\[
q_{00}(x)=q_{11}(x)=\frac v4,
\qquad
\frac v4\le q_{01}(x),q_{10}(x)\le1-\frac{3v}{4},
\qquad
\sum_{u_0,u_1\in\{0,1\}}q_{u_0,u_1}(x)=1.
\]
All twelve latent masses defining every \(P_{j,\ell}\in\mathcal W_n\) are nonnegative and sum to one; the always-taker and never-taker masses are strictly positive everywhere, and the complier masses are strictly positive wherever \(w_{j,\ell}>0\).  No H\"older restriction on these latent cells is used.  The observable one-sided extensions are
\[
\mu_{D,P_{j,\ell}}^{-}(x)=\frac{1-w_{j,\ell}(x)}2,
\qquad
\mu_{D,P_{j,\ell}}^{+}(x)=\frac{1+w_{j,\ell}(x)}2,
\]
\[
\mu_{Y,P_{j,\ell}}^{-}(x)=\frac12-\frac{g_\star(x)w_{j,\ell}(x)}2,
\qquad
\mu_{Y,P_{j,\ell}}^{+}(x)=\frac12+\frac{g_\star(x)w_{j,\ell}(x)}2.
\]
They belong to the radius-\(M\), order-\(s\) H\"older ball.  The potential reduced forms are continuous at \(c_{j,\ell}\), and
\[
d_{P_{j,\ell}}=\kappa_n,
\qquad
\theta_{P_{j,+}}=a_\star=1-v,
\qquad
\theta_{P_{j,-}}=-a_\star=-(1-v).
\]
Writing \(p_{j,\ell}^{D,Y}\) and \(p_{j,\ell}^{D}\) for the conditional observable cell probabilities, one has
\[
m_{d,y}(x)\ge\frac v2,
\qquad
p_{j,\ell}^{D,Y}(x,d,y)
=p_{j,\ell}^{D}(x,d)m_{d,y}(x)
\ge\frac{7v}{32}>0.
\]
On either observable side,
\[
\operatorname{Var}_{P_{j,\ell}}(D\mid X=x)=\frac{1-w_{j,\ell}(x)^2}{4}>0,
\]
\[
\operatorname{Var}_{P_{j,\ell}}(Y\mid X=x)
=\frac{1-g_\star(x)^2}{4}
+g_\star(x)^2\operatorname{Var}_{P_{j,\ell}}(D\mid X=x)
=\frac{1-g_\star(x)^2w_{j,\ell}(x)^2}{4}
\ge\frac{63}{256}>v,
\]
and
\[
\left|\operatorname{corr}_{P_{j,\ell}}(Y,D\mid X=x)\right|
=\frac{|g_\star(x)|\sqrt{1-w_{j,\ell}(x)^2}}
       {\sqrt{1-g_\star(x)^2w_{j,\ell}(x)^2}}
\le |g_\star(x)|\le1-v.
\]
The formula remains valid when \(w_{j,\ell}(x)=0\); treatment variance is never degenerate in this witness.  Consequently \(\mathcal W_n\subset\mathcal P_n\).  Moreover, with \(R_\star\) as the outcome-invariant reference,
\[
\frac{dP_{j,\ell}}{dR_\star}
=1+(2D-1)\operatorname{sgn}(X-c_{j,\ell})w_{j,\ell}(X)
\quad R_\star\text{-almost surely},
\]
so every product likelihood ratio is measurable with respect to \(\sigma((X_i,D_i)_{i=1}^n)\).  Finally,
\[
\int_{-1}^{1}\psi(u)^2\,du=\frac{256}{315},
\qquad
E_{R_\star}\!\left[
\frac{dP_{j,\ell}}{dR_\star}
\frac{dP_{k,r}}{dR_\star}
\right]=1
\]
for distinct \((j,\ell)\ne(k,r)\), and, for
\(Q_{\ell,n}=J_n^{-1}\sum_{j=1}^{J_n}P_{j,\ell}^{\otimes n}\),
\[
\chi^2(Q_{\ell,n},R_\star^{\otimes n})
=\frac{\left(1+\frac{256}{315}\kappa_n^2h_n\right)^n-1}{J_n}.
\]
\end{proposition}
\begin{proof}
Write \(a=a_\star=1-v\) and \(b=1-v/2\).  We verify each model property and likelihood calculation used by the lower bound.

First, \(h_n\le1/32\) implies \((16h_n)^{-1}\ge2\), and hence
\[
J_n=\left\lfloor(16h_n)^{-1}\right\rfloor\ge2.
\tag{1}
\]
For \(1\le j\le J_n\), the location formulas in \(\mathrm{def:outcome\mbox{-}invariant\mbox{-}family}\) give
\[
c_{j,+}-h_n=\frac14+2(j-1)h_n\ge\frac14,
\qquad
c_{j,+}+h_n=\frac14+2jh_n\le\frac38,
\tag{2}
\]
and similarly
\[
[c_{j,-}-h_n,c_{j,-}+h_n]\subset[5/8,3/4].
\tag{3}
\]
Centers within either label are separated by \(2h_n\), and the two labels use the disjoint plateau intervals in (2)--(3).  Because \(\psi\) is supported on \([-1,1]\), distinct supports meet at most at endpoints.  Uniformity of \(X\) makes those endpoints null, and therefore
\[
w_{j,\ell}(x)w_{k,r}(x)=0
\quad\text{for Lebesgue-almost every }x
\quad\text{whenever }(j,\ell)\ne(k,r).
\tag{4}
\]

Next, \(B'(t)=-30t^2(1-t)^2\le0\), \(B(0)=1\), and \(B(1)=0\).  Thus \(\mathrm{def:common\mbox{-}outcome\mbox{-}profile}\) gives \(|g_\star(x)|\le a=1-v\).  By the corrected construction in \(\mathrm{def:common\mbox{-}outcome\mbox{-}table}\),
\[
q_{00}=q_{11}=\frac v4,
\qquad
q_{01}=\frac{1-v/2+g_\star}{2},
\qquad
q_{10}=\frac{1-v/2-g_\star}{2}.
\tag{5}
\]
The four entries sum to one.  Since \(|g_\star|\le1-v\),
\[
\frac v4
=\frac{1-v/2-(1-v)}2
\le q_{01},q_{10}
\le\frac{1-v/2+(1-v)}2
=1-\frac{3v}{4}.
\tag{6}
\]
Together with \(v>0\), (5)--(6) show that every common-table cell is strictly positive.

The quartic profile obeys \(0\le\psi\le1\).  Hence \(0\le w_{j,\ell}\le\kappa_n\le1/8\), and the latent masses
\[
\pi_{C,u_0,u_1,P_{j,\ell}}=w_{j,\ell}q_{u_0,u_1},
\qquad
\pi_{A,u_0,u_1,P_{j,\ell}}
=\pi_{N,u_0,u_1,P_{j,\ell}}
=\frac{1-w_{j,\ell}}2q_{u_0,u_1}
\tag{7}
\]
are nonnegative.  The always-taker and never-taker cells are bounded below by
\[
\frac{1-1/8}{2}\frac v4=\frac{7v}{64}>0,
\tag{8}
\]
and a complier cell is positive wherever \(w_{j,\ell}>0\).  Summing (7) over response types and the four potential-outcome cells gives
\[
\sum_{u_0,u_1}q_{u_0,u_1}
\left(w_{j,\ell}+\frac{1-w_{j,\ell}}2+
                    \frac{1-w_{j,\ell}}2\right)=1.
\tag{9}
\]
This proves validity of the latent law using positivity alone; no smoothness of an individual latent cell is imposed or used.

We next verify the observable and potential reduced forms defining \(\mathcal P_n\).  The response-type weights in (7) imply
\[
E[D(z)\mid X=x]=\frac{1-w_{j,\ell}(x)}2+z\,w_{j,\ell}(x).
\tag{10}
\]
Because the potential-outcome table is common to every response type, (5) gives
\[
E[Y(0)\mid X=x]=\frac{1-g_\star(x)}2,
\qquad
E[Y(1)\mid X=x]=\frac{1+g_\star(x)}2.
\tag{11}
\]
Combining (10)--(11) yields
\[
E[Y(D(z))\mid X=x]
=\frac12+g_\star(x)\left(z-\frac12\right)w_{j,\ell}(x).
\tag{12}
\]
Equations (10)--(12), consistency, and \(Z=\mathbf 1\{X\ge c_{j,\ell}\}\) give the four displayed observable extensions in the statement.

It remains to check their H\"older radii.  On \((-1,1)\),
\[
\psi'(u)=-4u(1-u^2),
\qquad
\psi''(u)=-4+12u^2.
\tag{13}
\]
Both \(\psi\) and \(\psi'\) join continuously to zero at \(\{-1,1\}\).  Consequently
\[
\lVert\psi'\rVert_\infty\le2,
\qquad
\operatorname{Lip}(\psi')\le8.
\tag{14}
\]
When \(s=1\), the identity \(h_n=32\kappa_n/M\) and (14) give
\[
\operatorname{Lip}(w_{j,\ell})
\le\frac{2\kappa_n}{h_n}=\frac{M}{16}.
\tag{15}
\]
When \(1<s\le2\), put \(\beta=s-1\).  The bounded, globally Lipschitz function \(\psi'\) has \(\beta\)-H\"older seminorm at most \(8\): for arguments within unit distance this follows from (14), and for arguments farther apart it follows from \(2\lVert\psi'\rVert_\infty\le4\).  Since \(h_n^s=32\kappa_n/M\) and \(h_n\le1\), scaling gives
\[
\lVert w_{j,\ell}'\rVert_\infty
\le\frac{2\kappa_n}{h_n}
=\frac{M}{16}h_n^{s-1}\le\frac{M}{16},
\qquad
[w_{j,\ell}']_{s-1}
\le\frac{8\kappa_n}{h_n^s}=\frac{M}{4}.
\tag{16}
\]
By (2)--(3), \(g_\star=a\) on the support of every \(+\) profile and \(g_\star=-a\) on the support of every \(-\) profile.  Since each profile vanishes off its support,
\[
g_\star w_{j,+}=a w_{j,+},
\qquad
g_\star w_{j,-}=-a w_{j,-}
\quad\text{on }[0,1].
\tag{17}
\]
Thus, for \(s=1\), the treatment and outcome extensions have Lipschitz constants at most \(M/32\) and \(aM/32\), respectively.  For \(1<s\le2\), their derivative suprema are bounded by the same quantities and their derivative H\"older seminorms are at most \(M/8\) and \(aM/8\), respectively.  Since \(a<1\), all four extensions have supremum norm at most \(1/2+\kappa_n/2\le9/16\).  Even under the convention that the radius bounds the sum of these norm components,
\[
\frac9{16}+\frac{M}{32}+\frac{M}{8}\le M
\tag{18}
\]
because \(M\ge2\); deleting the derivative-seminorm term when \(s=1\) only strengthens the inequality.  Hence the extensions satisfy \(\mathrm{ass:observable\mbox{-}holder\mbox{-}extensions}\).  Equations (10), (12), and (17) are continuous at \(c_{j,\ell}\), so \(\mathrm{ass:potential\mbox{-}reduced\mbox{-}form\mbox{-}continuity}\) also holds.

Since \(w_{j,\ell}(c_{j,\ell})=\kappa_n>0\), the observable treatment extensions give
\[
d_{P_{j,\ell}}=\kappa_n.
\tag{19}
\]
The outcome extensions and (17) give
\[
y_{P_{j,+}}=a\kappa_n,
\qquad
y_{P_{j,-}}=-a\kappa_n.
\tag{20}
\]
Division by the positive denominator in (19) is therefore legitimate, and the definition of \(\theta_P\) yields the causal labels \(1-v\) and \(-(1-v)\).  The cutoff is interior by (2)--(3), \(X\) is uniform by construction, and taking independent copies gives the required iid sample.

For nondegeneracy and outcome invariance, the common table makes the conditional law of \(Y\) given \((D=d,X=x)\) Bernoulli with success probability
\[
E[Y\mid D=d,X=x]=\frac12+g_\star(x)\left(d-\frac12\right).
\tag{21}
\]
Every success or failure probability in (21) is at least
\[
\frac{1-|g_\star(x)|}{2}\ge\frac v2.
\tag{22}
\]
On either side of the cutoff, the two conditional treatment probabilities are \((1+w_{j,\ell})/2\) and \((1-w_{j,\ell})/2\), in an order determined by the side, so each is at least \(7/16\).  Therefore the common-table factorization gives
\[
p_{j,\ell}^{D,Y}(x,d,y)
=p_{j,\ell}^{D}(x,d)m_{d,y}(x)
\ge\frac{7v}{32}>0.
\tag{23}
\]
This strict positivity permits density cancellation below; no comparison of the joint-cell floor with the model constant \(v\) is needed or asserted.

Let \(w=w_{j,\ell}(x)\) and \(g=g_\star(x)\).  On either side, \(D\mid X=x\) is Bernoulli with parameter \((1+w)/2\) or \((1-w)/2\), and hence
\[
V_D:=\operatorname{Var}(D\mid X=x)=\frac{1-w^2}{4}\ge\frac{63}{256}>0.
\tag{24}
\]
The conditional variance identity and (21) give
\[
\operatorname{Var}(Y\mid X=x)
=\frac{1-g^2}{4}+g^2V_D
=\frac{1-g^2w^2}{4}
\ge\frac{63}{256}>\frac1{64}>v,
\tag{25}
\]
where the first lower bound uses \(|g|\le1\) and \(w\le1/8\), and the last uses \(0<v<1/64\).  Likewise,
\[
\operatorname{Cov}(Y,D\mid X=x)=gV_D.
\tag{26}
\]
Because (24) is strictly positive, (24)--(26) imply the exact formula
\[
|\operatorname{corr}(Y,D\mid X=x)|
=\frac{|g|\sqrt{1-w^2}}{\sqrt{1-g^2w^2}}
\le |g|\le1-v.
\tag{27}
\]
Indeed, \(g^2\le1\) gives \(1-w^2\le1-g^2w^2\).  Formula (27) also applies at \(w=0\), where its square-root ratio equals one.  No treatment-degenerate point occurs in this witness by (24); if such a point occurred in a different law, (26) would give zero covariance and the convention in \(\mathrm{ass:observable\mbox{-}correlation\mbox{-}separation}\) would apply without forming a correlation coefficient.  Equations (2)--(27) verify every member property in \(\mathrm{def:model\mbox{-}class}\), so
\[
\mathcal W_n\subset\mathcal P_n.
\tag{28}
\]

Finally, under \(R_\star\), \(D\) is conditionally fair given \(X\), while (21) is the same strictly positive conditional outcome margin as under every \(P_{j,\ell}\).  Cancelling that margin in the full observable density ratio gives, outside the null event \(X=c_{j,\ell}\),
\[
L_{j,\ell}:=\frac{dP_{j,\ell}}{dR_\star}
=1+(2D-1)\operatorname{sgn}(X-c_{j,\ell})w_{j,\ell}(X).
\tag{29}
\]
Thus the one-unit and product likelihood ratios are measurable with respect to the corresponding \((X,D)\) sigma-fields.  Direct integration gives
\[
\int_{-1}^{1}\psi(u)^2\,du
=2\left(1-\frac43+\frac65-\frac47+\frac19\right)
=\frac{256}{315}.
\tag{30}
\]
Conditioning on \(X\) under \(R_\star\) in (29), and using \(E_{R_\star}[2D-1\mid X]=0\), yields
\[
E_{R_\star}[L_{j,\ell}L_{k,r}]
=1+\int_0^1
\operatorname{sgn}(x-c_{j,\ell})
\operatorname{sgn}(x-c_{k,r})
w_{j,\ell}(x)w_{k,r}(x)\,dx.
\tag{31}
\]
For distinct indices, (4) makes (31) equal to one.  For equal indices, the substitution \(u=(x-c_{j,\ell})/h_n\) and (30) give
\[
E_{R_\star}[L_{j,\ell}^2]
=1+\frac{256}{315}\kappa_n^2h_n.
\tag{32}
\]
Independence raises (31)--(32) to the \(n\)-th power.  Expanding the square of the equally weighted mixture likelihood ratio gives \(J_n\) diagonal terms and \(J_n(J_n-1)\) off-diagonal terms.  Subtracting one therefore yields
\[
\chi^2(Q_{\ell,n},R_\star^{\otimes n})
=\frac{\left(1+\frac{256}{315}\kappa_n^2h_n\right)^n-1}{J_n},
\]
which completes the proof.
\end{proof}
\par\smallskip\noindent \textit{Justification.} The explicit positivity, observable smoothness, exact variance and correlation, factorization, and second-moment calculations make outcome-assisted localization fail on a legal parameter-exact high-separation causal subexperiment. \textit{Gap.} Aligned multivariate change-point analyses do not provide a monotone fuzzy-RD latent construction whose outcome margin cancels exactly while its two LATE labels are separated by \(2(1-v)\). \textit{Consumer.} The finite-sample lower theorem uses the family to transfer honesty between the opposite labels \(\pm(1-v)\) without allowing the outcome channel to eliminate the location-entropy barrier.

\begin{theorem}[thm:finite-sample-location-barrier-explicit-elbow]\label{thm:finite-sample-location-barrier-explicit-elbow}
With \(a_\star=1-v\), put
\[
C_\psi=\frac{256}{315},
\qquad
C_h=\left(\frac{32}{M}\right)^{1/s},
\qquad
\lambda_w=(sC_\psi C_h)^{-1}
=\frac{315}{256s}\left(\frac{M}{32}\right)^{1/s}.
\]
If \(h_n\le1/32\), then the high-separation outcome-invariant witness gives the finite-sample bound
\[
L_n(\kappa_n)
\ge2(1-v)\left[
1-2\alpha-
\sqrt{\frac{\left(1+C_\psi\kappa_n^2h_n\right)^n-1}{J_n}}
\right]_+.
\]
Along every admissible triangular sequence,
\[
\frac1s\log(1/\kappa_n)-C_\psi C_hT_n\longrightarrow+\infty
\quad\Longrightarrow\quad
\liminf_{n\to\infty}L_n(\kappa_n)
\ge2(1-v)(1-2\alpha).
\]
In particular,
\[
\limsup_{n\to\infty}
\frac{T_n}{\log(1/\kappa_n)}<\lambda_w
\quad\Longrightarrow\quad
\liminf_{n\to\infty}L_n(\kappa_n)
\ge2(1-v)(1-2\alpha).
\]
For \(s=1\), \(\lambda_w=315M/8192\).  The sufficient regime contains sequences on which the known-cutoff minimax length shrinks to zero: for every \(0<\gamma<(2+1/s)^{-1}\), the eventual choice \(\kappa_n=\rho_n(\log n)^\gamma\) satisfies the first displayed limit and \(L_n^{\mathrm{known}}(\kappa_n)\to0\).  The constant \(\lambda_w\) is an explicit one-sided witness threshold; this theorem does not claim that it is the exact critical frontier.
\end{theorem}
\begin{proof}
Let \(I_n\) be an arbitrary measurable connected interval rule satisfying the uniform coverage condition in \(\mathrm{def:minimax\mbox{-}length}\).  By \(\mathrm{prop:outcome\mbox{-}invariant\mbox{-}witness\mbox{-}high\mbox{-}separation}\), each component of \(Q_{+,n}\) is in \(\mathcal P_n\) and has causal target \(a_\star=1-v\), while each component of \(Q_{-,n}\) is in \(\mathcal P_n\) and has target \(-a_\star=-(1-v)\).  Averaging the componentwise honesty inequalities gives
\[
Q_{+,n}\{a_\star\in I_n\}\ge1-\alpha,
\qquad
Q_{-,n}\{-a_\star\in I_n\}\ge1-\alpha.
\tag{33}
\]
For probability measures \(Q\ll R\), Cauchy--Schwarz applied to the density ratio proves
\[
\operatorname{TV}(Q,R)
=\frac12E_R\left|\frac{dQ}{dR}-1\right|
\le\frac12\sqrt{E_R\left(\frac{dQ}{dR}-1\right)^2}
=\frac12\sqrt{\chi^2(Q,R)}.
\tag{34}
\]
The triangle inequality for total variation, (34), and the equal chi-square identities for the two labels yield
\[
\operatorname{TV}(Q_{+,n},Q_{-,n})
\le\operatorname{TV}(Q_{+,n},R_\star^{\otimes n})
   +\operatorname{TV}(Q_{-,n},R_\star^{\otimes n})
\le
\sqrt{\frac{\left(1+C_\psi\kappa_n^2h_n\right)^n-1}{J_n}}.
\tag{35}
\]
Transfer the second event in (33) from \(Q_{-,n}\) to \(Q_{+,n}\) by the definition of total variation, and intersect it with the first event in (33).  The union bound and (35) give
\[
Q_{+,n}\{-a_\star,a_\star\in I_n\}
\ge1-2\alpha-
\sqrt{\frac{\left(1+C_\psi\kappa_n^2h_n\right)^n-1}{J_n}}.
\tag{36}
\]
A connected subset of \(\Theta=[-1,1]\) containing both labels has Lebesgue length at least \(2a_\star=2(1-v)\).  Since \(Q_{+,n}\) is the uniform mixture of in-class product laws,
\[
\sup_{P\in\mathcal P_n}E_P[|I_n|]
\ge E_{Q_{+,n}}[|I_n|]
\ge2(1-v)\left[
1-2\alpha-
\sqrt{\frac{\left(1+C_\psi\kappa_n^2h_n\right)^n-1}{J_n}}
\right]_+.
\tag{37}
\]
Taking the infimum over all honest rules proves the finite-sample assertion.

We now evaluate (37).  The definitions of \(h_n\), \(C_h\), and \(T_n\) imply
\[
h_n=C_h\kappa_n^{1/s},
\qquad
n\kappa_n^2h_n=C_hT_n.
\tag{38}
\]
Whenever \(h_n\le1/32\), the number \(x=(16h_n)^{-1}\) is at least two.  Since \(\lfloor x\rfloor\ge x/2\) for \(x\ge2\),
\[
J_n\ge\frac1{32h_n}.
\tag{39}
\]
Using \(1+x\le e^x\), equations (38)--(39) bound the chi-square term in (37) by
\[
\frac{\left(1+C_\psi\kappa_n^2h_n\right)^n-1}{J_n}
\le32h_n\exp\{C_\psi n\kappa_n^2h_n\}
=32C_h\kappa_n^{1/s}\exp\{C_\psi C_hT_n\}.
\tag{40}
\]
The logarithm of the final expression is
\[
\log(32C_h)-\frac1s\log(1/\kappa_n)+C_\psi C_hT_n.
\tag{41}
\]
Suppose
\[
\frac1s\log(1/\kappa_n)-C_\psi C_hT_n\longrightarrow+\infty.
\tag{42}
\]
Because \(T_n>0\), condition (42) implies \(\log(1/\kappa_n)\to\infty\).  Hence \(\kappa_n\to0\), so \(h_n\to0\) and (39)--(41) apply eventually.  Equation (41) tends to \(-\infty\), and thus the square-root term in (37) tends to zero.  We conclude that
\[
\liminf_{n\to\infty}L_n(\kappa_n)
\ge2(1-v)(1-2\alpha).
\tag{43}
\]

For the ratio corollary, suppose
\[
\limsup_n\frac{T_n}{\log(1/\kappa_n)}<\lambda_w.
\tag{44}
\]
Condition (44) forces \(\kappa_n\to0\).  Otherwise, along a subsequence bounded away from zero, \(T_n=n\kappa_n^{2+1/s}\) would diverge at least linearly while \(\log(1/\kappa_n)\) remained bounded, contradicting (44).  Choose \(\varepsilon>0\) so that eventually
\[
\frac{T_n}{\log(1/\kappa_n)}\le\lambda_w-\varepsilon.
\]
Since \(C_\psi C_h\lambda_w=1/s\),
\[
\frac1s\log(1/\kappa_n)-C_\psi C_hT_n
\ge C_\psi C_h\varepsilon\log(1/\kappa_n)
\longrightarrow+\infty.
\tag{45}
\]
Thus (44) implies (42), and (43) proves the stated corollary.  Substitution of \(s=1\) gives
\[
\lambda_w=\frac{315}{256}\frac{M}{32}=\frac{315M}{8192}.
\]

Finally, fix \(0<\gamma<(2+1/s)^{-1}\) and, for all sufficiently large \(n\), set
\[
\kappa_n=\rho_n(\log n)^\gamma
=n^{-s/(2s+1)}(\log n)^\gamma.
\tag{46}
\]
This sequence is positive and eventually at most \(1/8\).  Direct calculation gives
\[
T_n=n\kappa_n^{2+1/s}
=(\log n)^{\gamma(2+1/s)}=o(\log n),
\tag{47}
\]
while
\[
\log(1/\kappa_n)
=\frac{s}{2s+1}\log n-\gamma\log\log n
\asymp\log n.
\tag{48}
\]
Equations (47)--(48) imply (42).  Moreover,
\[
\frac{\rho_n}{\kappa_n}=(\log n)^{-\gamma}\longrightarrow0.
\]
The upper inequality in \(\mathrm{thm:known\mbox{-}cutoff\mbox{-}benchmark}\) therefore gives \(L_n^{\mathrm{known}}(\kappa_n)\to0\), whereas (43) remains in force.  This proves the oracle-shrinking comparison.  No converse or attaining construction is supplied at or above \(\lambda_w\), so the conclusion remains the asserted one-sided witness threshold.
\end{proof}
\par\smallskip\noindent \textit{Justification.} The bound gives an explicit finite-sample connected-length certificate and evaluates the largest one-sided signal-to-entropy threshold delivered by this exact witness calculation. \textit{Gap.} Known-cutoff weak-instrument inference and shrinking-jump cutoff-location rates do not imply the \(2(1-v)\) connected LATE-length bound or the explicit threshold \(\lambda_w\) under a jointly unknown and vanishing first stage. \textit{Consumer.} Analysts can compute \(\lambda_w\) and identify sequences on which an administratively undocumented cutoff makes every honest connected LATE interval nonshrinking, despite a shrinking known-cutoff benchmark.

\section{Interpretation}

The ratio \(T_n/\log(1/\kappa_n)\) compares local jump information with the number of disjoint candidate cutoffs.  Below the explicit witness threshold \(\lambda_w=(sC_\psi C_h)^{-1}\), the two mixtures with causal labels \(\pm(1-v)\) merge in total variation, so honesty and connectedness force asymptotic expected length at least \(2(1-v)(1-2\alpha)\), even along sequences for which known-cutoff intervals shrink.  This sufficient converse does not locate the exact critical transition.  At \(s=1\), strong signal permits the proved rectangular scan and oracle-order inversion.  For \(s>1\), the frozen affine profile fails at the population level, so a different criterion or a negative frontier theorem is necessary.

\section*{Technical internal limitation (diagnostic only)}

The lower construction evaluates only one side of the critical experiment.  At \(T_n/\log(1/\kappa_n)\to\lambda_w\), the first-order coordinate does not determine the witness mixture limit because the second-order centering \(C_\psi C_hT_n-s^{-1}\log(1/\kappa_n)\) remains; above \(\lambda_w\), this chi-square argument becomes uninformative.  The higher-smoothness profile branch remains a total measurable map but is not a valid uniform localizer on the stated class: a legal population score favors a displaced split, and maximal inequalities cannot repair that deterministic failure.  Neither limitation affects the finite-sample \(2(1-v)\) lower bound or the proved Lipschitz-face upper theorem, but the exact critical law, a replacement higher-smoothness estimator, and a matching slope-kink lower pair remain open.

\section{Honest scope}

The Wald interpretation, known-cutoff benchmark, strengthened finite-sample location-mixture barrier, and Lipschitz-face oracle attainment retain their stated scope.  The explicit \(\lambda_w\) condition is a sufficient one-sided witness threshold and is not advertised as the exact signal-to-entropy boundary.  The coefficient \(2(1-v)\) comes from the legal common-outcome witness with \(a_\star=1-v\) and \(b=1-v/2\); no latent-cell H{\"o}lder restriction is imposed or claimed, and no observable joint-cell floor larger than \(v\) is required or asserted.  For \(s>1\), the paper proves only that the frozen common-versus-piecewise moving-window profile fails its proposed uniform localization inequality on the unchanged class; it does not conclude that every estimator fails or that the candidate detectable-region bracket is false.  No cross-side derivative matching, adaptation guarantee, disconnected-set result, or matched critical-window constant is added.  Noack--Rothe-style refers only to the observable-regression and nondegeneracy template; equality with, or inclusion in, their published parameter space is not claimed.

\begin{thebibliography}{99}

  \bibitem{HahnToddVanDerKlaauw2001} Hahn, Jinyong, Petra Todd, and Wilbert van der Klaauw (2001), ``Identification and Estimation of Treatment Effects with a Regression-Discontinuity Design,'' Econometrica 69(1), 201--209.

  \bibitem{MullerSong1997} M{\"u}ller, Hans-Georg, and Kai-Sheng Song (1997), ``Two-stage change-point estimators in smooth regression models,'' Statistics and Probability Letters 34(4), 323--335, doi:10.1016/S0167-7152(96)00197-6.

  \bibitem{PorterYu2015} Porter, Jack, and Ping Yu (2015), ``Regression discontinuity designs with unknown discontinuity points: Testing and estimation,'' Journal of Econometrics 189(1), 132--147, doi:10.1016/j.jeconom.2015.06.002.

  \bibitem{ArmstrongKolesar2018} Armstrong, Timothy B., and Michal Koles{\'a}r (2018), ``Optimal Inference in a Class of Regression Models,'' Econometrica 86(2), 655--683, doi:10.3982/ECTA14434, arXiv:1511.06028.

  \bibitem{FrickMunkSieling2014} Frick, Klaus, Axel Munk, and Hannes Sieling (2014), ``Multiscale change point inference,'' Journal of the Royal Statistical Society: Series B 76(3), 495--580, doi:10.1111/rssb.12047, arXiv:1301.7212.

  \bibitem{ChernozhukovChetverikovKato2014} Chernozhukov, Victor, Denis Chetverikov, and Kengo Kato (2014), ``Gaussian Approximation of Suprema of Empirical Processes,'' Annals of Statistics 42(4), 1564--1597, arXiv:1212.6885.

  \bibitem{AndrewsKitagawaMcCloskey2021} Andrews, Isaiah, Toru Kitagawa, and Adam McCloskey (2021), ``Inference after estimation of breaks,'' Journal of Econometrics 224(1), 39--59, doi:10.1016/j.jeconom.2020.07.036.

  \bibitem{NoackRothe2024} Noack, Claudia, and Christoph Rothe (2024), ``Bias-Aware Inference in Fuzzy Regression Discontinuity Designs,'' Econometrica 92(3), 687--711, doi:10.3982/ECTA19466, arXiv:1906.04631.

  \bibitem{Tuvaandorj2020} Tuvaandorj, Purevdorj (2020), ``Regression discontinuity designs, white noise models, and minimax,'' Journal of Econometrics 218(2), 587--608, doi:10.1016/j.jeconom.2020.04.030.

  \bibitem{MadridPadillaYuWangRinaldo2022} Madrid Padilla, Oscar Hernan, Yi Yu, Daren Wang, and Alessandro Rinaldo (2022), ``Optimal nonparametric multivariate change point detection and localization,'' IEEE Transactions on Information Theory, arXiv:1910.13289.

  \bibitem{KowalskaVanDeWielVanDerPas2025} Kowalska, Julia M., Mark A. van de Wiel, and St{\'e}phanie van der Pas (2025 version), ``Bayesian Regression Discontinuity Design with Unknown Cutoff,'' arXiv:2406.11585.

  \bibitem{AraiOtsuSeo2026} Arai, Yoichi, Taisuke Otsu, and Myung Hwan Seo (2026), ``Causal Inference on Regression Discontinuity Designs by High-Dimensional Methods,'' KEO Discussion Paper 200.

\end{thebibliography}

\end{document}